%% 3-8 かっこをはずす：展開と平方完成（第12回：新しく書いた。組 V3-4）  % noruby
%% 第5巻（5-2・5-4・5-6・5-7・5-8）が指す形：(a+b)^2、(a-b)^2、(1+h)^2、(12-2x)^2、(a-b)(a+b)、「全部の組み合わせをかけてたす」、立方体の色分け。  % noruby
%% 第11回の決定：恒等式・係数を比べる・数を代入して係数を決める（9-7 の部分分数分解・ラプラス変換の入口）。2通り以上で説明する。  % noruby
\makeatletter
% 斜めから見た箱（5-4 の \bkvebox と同じ形・同じぬり方）  \bkthribox{x0}{y0}{z0}{dx}{dy}{dz}{色}  % noruby
\newcommand\bkthribox[7]{%
  \pgfmathsetmacro\bkthriX{#1+#4}\pgfmathsetmacro\bkthriY{#2+#5}\pgfmathsetmacro\bkthriZ{#3+#6}%
  \filldraw[fill=#7!42, draw=#7!60!black, line width=0.5pt, line join=round] (#1,#2,#3) -- (\bkthriX,#2,#3) -- (\bkthriX,#2,\bkthriZ) -- (#1,#2,\bkthriZ) -- cycle;
  \filldraw[fill=#7!18, draw=#7!60!black, line width=0.5pt, line join=round] (#1,#2,\bkthriZ) -- (\bkthriX,#2,\bkthriZ) -- (\bkthriX,\bkthriY,\bkthriZ) -- (#1,\bkthriY,\bkthriZ) -- cycle;
  \filldraw[fill=#7!68, draw=#7!60!black, line width=0.5pt, line join=round] (\bkthriX,#2,#3) -- (\bkthriX,\bkthriY,#3) -- (\bkthriX,\bkthriY,\bkthriZ) -- (\bkthriX,#2,\bkthriZ) -- cycle;}
% 1辺 a の立方体と、1辺を a+b にしたときにふえる7つの部分（5-4 の \bkvecube と同じ色分け）  % noruby
%   \bkthricube{a}{b}{右のすき間}{前のすき間}{上のすき間}  青：もとの立方体、オレンジ：板3まい、緑：角柱3本、紫：小さな立方体  % noruby
\newcommand\bkthricube[5]{%
  \bkthribox{0}{0}{0}{#1}{#1}{#1}{bkBlue}%
  \bkthribox{#1+#3}{0}{0}{#2}{#1}{#1}{bkOrange}%
  \bkthribox{0}{0}{#1+#5}{#1}{#1}{#2}{bkOrange}%
  \bkthribox{0}{-#4-#2}{0}{#1}{#2}{#1}{bkOrange}%
  \bkthribox{#1+#3}{-#4-#2}{0}{#2}{#2}{#1}{bkGreen}%
  \bkthribox{0}{-#4-#2}{#1+#5}{#1}{#2}{#2}{bkGreen}%
  \bkthribox{#1+#3}{0}{#1+#5}{#2}{#1}{#2}{bkGreen}%
  \bkthribox{#1+#3}{-#4-#2}{#1+#5}{#2}{#2}{#2}{bkPurple}}
% 天びん（bk-pict1 の \bkbalance）の皿に、数を書いたおもりをのせる  \bkthribal{(x,y)}{かたむき}{左の数}{右の数}  % noruby
\newcommand\bkthribal[4]{%
  \bkbalance{#1}{#2}%
  \node[draw=bkBlue!70!black, fill=bkBlue!15, rounded corners=0.6mm, font=\sffamily\small, inner sep=1.2pt, minimum width=9mm, anchor=south] at (balL) {#3};
  \node[draw=bkOrange!80!black, fill=bkOrange!18, rounded corners=0.6mm, font=\sffamily\small, inner sep=1.2pt, minimum width=9mm, anchor=south] at (balR) {#4};}
% 式を、まん中に1行で置く  \bkthriline{…}  % noruby
\newcommand\bkthriline[1]{\par\addvspace{3pt}{\centering #1\par}\addvspace{3pt}}
% 説明の段の見出し（丸数字＋太字の1行）  \bkthristep{番号}{見出し}  % noruby
\newcommand\bkthristep[2]{\par\addvspace{6pt}\Needspace{4\baselineskip}\noindent
  {\bkcircled[bkOrange]{#1}\enspace\sffamily\bfseries\color{bkOrange!75!black}#2}\par\nobreak\vspace{2pt}}
% 本文の中の小見出し（それだけで1行）  \bkthrihead{見出し}  % noruby
\newcommand\bkthrihead[1]{\par\addvspace{4pt}\Needspace{3\baselineskip}\noindent{\sffamily\bfseries #1}\par\nobreak\noindent}
% 枠でかこんだ大事な式  \bkthribox{…} とまちがえないように名前を分ける  \bkthriwaku{…}  % noruby
\newcommand\bkthriwaku[1]{\par\addvspace{4pt}{\centering\setlength\fboxsep{4pt}\fcolorbox{bkOrange}{white}{#1}\par}\addvspace{4pt}}
\makeatother

\bkchapter{3}{8}{かっこをはずす：\\\R{展開}{てん|かい}と\R{平方完成}{へい|ほう|かん|せい}}{\bklevelrange{2}{3}}{かっこをはずす：\R{展開}{てん|かい}と\R{平方完成}{へい|ほう|かん|せい}}
\bkmichishirube{2-11,2-12,2-13,3-6}{3-9,3-10,3-16,5-4,9-7}

\begin{goal}
\begin{itemize}[label={\color{\bkTT}●}, leftmargin=1.6\zw, itemsep=0pt, topsep=1pt]
\item \bklevel{2}\enspace 「\R{全部}{ぜん|ぶ}の\R{組}{く}み\R{合}{あ}わせをかけてたす」で、かっこのある\R{式}{しき}を\R{展開}{てん|かい}できる。\R{面積}{めん|せき}の\R{図}{ず}でわけを\R{説明}{せつ|めい}できる。
\item \bklevel{3}\enspace \R{恒等式}{こう|とう|しき}と\R{方程式}{ほう|てい|しき}を\R{見分}{み|わ}け、\R{係数}{けい|すう}を\R{比}{くら}べたり\R{数}{かず}を\R{代入}{だい|にゅう}したりして、わからない\R{係数}{けい|すう}を\R{決}{き}められる。
\item \bklevel{3}\enspace \R{平方完成}{へい|ほう|かん|せい}をして、\R{式}{しき}のいちばん\R{小}{ちい}さい\R{値}{あたい}・いちばん\R{大}{おお}きい\R{値}{あたい}を\R{読}{よ}みとれる。
\end{itemize}
\end{goal}

\begin{hatena}
\noindent\begin{minipage}[c]{0.62\linewidth}
\R{真上}{ま|うえ}に\qty{20}{m/s}で\R{投}{な}げ\R{上}{あ}げたボールの、$t$ \R{秒後}{びょう|ご}の\R{高}{たか}さは $h=20t-5t^2$（m）です（わけは5-5で）。$t=0$〜4では $h=0$、15、20、15、0 なので、2\R{秒後}{びょう|ご}がいちばん\R{高}{たか}そう。でも、1.9\R{秒}{びょう}や2.1\R{秒}{びょう}のほうが\R{高}{たか}いかも？

この\R{式}{しき}は $h=20-5(t-2)^2$ とも\R{書}{か}けて、この\R{形}{かたち}ならいちばん\R{高}{たか}い\R{点}{てん}がすぐ\R{見}{み}えるそうです。\R{本当}{ほん|とう}に\R{同}{おな}じ\R{式}{しき}？　なぜ\R{見}{み}えるの？
\end{minipage}\hfill
\begin{minipage}[c]{0.35\linewidth}\centering
\begin{tikzpicture}[x=1mm, y=0.85mm]
  \fill[bkGreen!25] (-4,-2) rectangle (33,0);
  \draw[bkGray, line width=0.6pt] (0,0) -- (0,42);
  \foreach \hh in {0,10,20} {\draw[bkGray, line width=0.5pt] (-1,{\hh*2}) -- (1,{\hh*2});
    \node[font=\sffamily\scriptsize, anchor=east] at (-1.5,{\hh*2}) {\hh};}
  \node[font=\sffamily\scriptsize, anchor=south] at (0,43) {$h$（m）};
  \foreach \tt/\hh/\xx in {0/0/8, 1/15/8, 2/20/16, 3/15/24, 4/0/24} {
    \shade[ball color=bkOrange] (\xx,{\hh*2+1.8}) circle (1.8);
    \node[font=\sffamily\scriptsize, anchor=west] at (\xx+2.2,{\hh*2+1.8}) {\tt\R{秒}{びょう}};}
  \draw[bk arrow, bkBlue, line width=0.8pt] (5,4) -- (5,30);
  \draw[bk arrow, bkRed, line width=0.8pt] (21,30) -- (21,4);
  \node[font=\sffamily\bfseries\large, text=bkOrange!80!black] at (16,37) {？};
\end{tikzpicture}
\end{minipage}
\end{hatena}

この\R{章}{しょう}では、かっこのある\R{式}{しき}のかっこをはずす\R{計算}{けい|さん}（\R{展開}{てん|かい}）と、その\R{反対}{はん|たい}に\R{式}{しき}をかっこにまとめる\R{計算}{けい|さん}を\R{学}{まな}びます。どちらも、\R{長方形}{ちょう|ほう|けい}や\R{正方形}{せい|ほう|けい}の\R{面積}{めん|せき}で\R{考}{かんが}えると、わけがよくわかります。\R{式}{しき}の\R{形}{かたち}を\R{自由}{じ|ゆう}に\R{変}{か}えられるようになると、\R{同}{おな}じ\R{式}{しき}から、ちがうことが\R{読}{よ}みとれるようになります。

%% ---------------- やってみよう・思い出そう ----------------
\begin{yattemiyou}[13×13を、4つに\R{分}{わ}けて\R{計算}{けい|さん}しよう]
\noindent\begin{minipage}[c]{0.56\linewidth}
\begin{enumerate}[label={\bkcircled[bkGreen]{\arabic*}}, leftmargin=2.4\zw, itemsep=1pt, topsep=0pt]
\item \R{方眼紙}{ほう|がん|し}に、1\R{辺}{ぺん}が13マスの\R{正方形}{せい|ほう|けい}をかく。
\item はしから10マスの\R{所}{ところ}に、たてと\R{横}{よこ}に1\R{本}{ぽん}ずつ\R{線}{せん}を\R{引}{ひ}いて、4つに\R{分}{わ}ける。
\item 4つの\R{部分}{ぶ|ぶん}のマスの\R{数}{かず}を、かけ\R{算}{ざん}で\R{求}{もと}めて、たそう。
\item $13\times13$ や、$10\times10+3\times3$ とくらべよう。
\end{enumerate}
\end{minipage}\hfill
\begin{minipage}[c]{0.4\linewidth}\centering
\begin{tikzpicture}[x=3.2mm, y=3.2mm]
  \fill[bkBlue!14] (0,3) rectangle (10,13);
  \fill[bkOrange!25] (10,3) rectangle (13,13);
  \fill[bkOrange!25] (0,0) rectangle (10,3);
  \fill[bkRed!35] (10,0) rectangle (13,3);
  \draw[bkLine, line width=0.3pt] (0,0) grid (13,13);
  \draw[bkInk, line width=0.8pt] (0,0) rectangle (13,13);
  \draw[bkInk, line width=0.8pt] (10,0) -- (10,13) (0,3) -- (13,3);
  \node[font=\sffamily\footnotesize, anchor=south] at (5,13.2) {10};
  \node[font=\sffamily\footnotesize, anchor=south] at (11.5,13.2) {3};
  \node[font=\sffamily\footnotesize, anchor=east] at (-0.2,8) {10};
  \node[font=\sffamily\footnotesize, anchor=east] at (-0.2,1.5) {3};
\end{tikzpicture}
\end{minipage}
\end{yattemiyou}

\noindent \R{青}{あお}い\R{部分}{ぶ|ぶん}は $10\times10=100$、オレンジの2つは $10\times3=30$ ずつ、\R{赤}{あか}い\R{部分}{ぶ|ぶん}は $3\times3=9$ マスです。たすと $100+30+30+9=169$ で、$13\times13=169$ とぴったり\R{合}{あ}います。いっぽう、$10\times10+3\times3=109$ では、オレンジの $30+30=60$ マスがたりません。「たしてから2\R{乗}{じょう}」と「2\R{乗}{じょう}してからたす」はちがうのです。

\begin{omoidasou}
\bkrecall{2-11}{\R{分配法則}{ぶん|ぱい|ほう|そく}：$(a+b)\times c=a\times c+b\times c$。\R{長方形}{ちょう|ほう|けい}の\R{面積}{めん|せき}を2\R{通}{とお}りに\R{表}{あらわ}した\R{式}{しき}。}
\bkrecall{2-12}{×を\R{省}{はぶ}く。\R{同類項}{どう|るい|こう}をまとめる（$3x+2x=5x$）。$3(x+2)=3x+6$。$x^2\times x^3=x^5$（2-7）。}
\bkrecall{2-13}{\R{式}{しき}をかっこごと\R{代入}{だい|にゅう}する：$x^2$ の $x$ に $a+h$ を\R{入}{い}れると $(a+h)^2$。}
\bkrecall{2-16}{\R{方程式}{ほう|てい|しき}は、\R{文字}{も|じ}に\R{入}{い}れる\R{数}{かず}によって、\R{成}{な}り\R{立}{た}ったり\R{成}{な}り\R{立}{た}たなかったりする\R{等式}{とう|しき}。}
\bkrecall{3-6}{\R{連立方程式}{れん|りつ|ほう|てい|しき}は、2つの\R{式}{しき}をたしたりひいたりして、\R{文字}{も|じ}を1つ\R{消}{け}して\R{解}{と}く（\R{加減法}{か|げん|ほう}）。}
\end{omoidasou}

%% ---------------- 数学の道具 ----------------
\bksection{doguu}{\R{数学}{すう|がく}の\R{道具}{どう|ぐ}}

\bksubsection{1}{かけ\R{算}{ざん}だけの\R{式}{しき}：\R{単項式}{たん|こう|しき}と\R{次数}{じ|すう}}

\leavevmode\label{bk:3-8:tankou}2-12では、$3\times a\times a\times b$ を $3a^2b$ と\R{短}{みじか}く\R{書}{か}く\R{約束}{やく|そく}をしました。$3a^2b$ や $5x$、$-2x^2$、7 のように、\R{数}{かず}と\R{文字}{も|じ}のかけ\R{算}{ざん}だけでできている\R{式}{しき}を\textbf{\R{単項式}{たん|こう|しき}}といいます。「\R{項}{こう}が1つだけの\R{式}{しき}」という\R{意味}{い|み}です。

$3x^2-2x+5$ のように、いくつかの\R{項}{こう}（＋でつながった1つ1つ、2-12）の\R{和}{わ}になっている\R{式}{しき}を\textbf{\R{多項式}{た|こう|しき}}といいます。ひき\R{算}{ざん}は「マイナスの\R{項}{こう}をたす」と\R{考}{かんが}えるので、$3x^2-2x+5=3x^2+(-2x)+5$ の\R{項}{こう}は $3x^2$、$-2x$、5 の3つです。

\bkthrihead{\R{次数}{じ|すう}：\R{文字}{も|じ}をいくつかけたか}
\R{単項式}{たん|こう|しき}で、かけ\R{合}{あ}わせてある\R{文字}{も|じ}の\R{個数}{こ|すう}を、その\R{単項式}{たん|こう|しき}の\textbf{\R{次数}{じ|すう}}といいます。\bkidx{じすう}{\R{次数}{じ|すう}}

\begin{bkcenter}
\begin{bktable}{colspec={Q[c,m,wd=20mm] Q[c,m,wd=40mm] Q[c,m,wd=26mm] Q[c,m,wd=16mm]}, rowsep=1.5pt, cells={font=\small}}
\R{単項式}{たん|こう|しき} & ×を\R{書}{か}くと & かけた\R{文字}{も|じ}の\R{数}{かず} & \R{次数}{じ|すう} \\
$5x$ & $5\times x$ & 1こ & 1 \\
$-2x^2$ & $(-2)\times x\times x$ & 2こ & 2 \\
$3a^2b$ & $3\times a\times a\times b$ & 3こ & 3 \\
7 & 7 & 0こ & 0 \\
\end{bktable}
\bktabcaption{\R{次数}{じ|すう}は、×を\R{全部}{ぜん|ぶ}\R{書}{か}いて\R{文字}{も|じ}を\R{数}{かぞ}えるとわかる。}
\end{bkcenter}

\R{多項式}{た|こう|しき}の\R{次数}{じ|すう}は、いちばん\R{次数}{じ|すう}の\R{大}{おお}きい\R{項}{こう}の\R{次数}{じ|すう}です。いちばん\R{大}{おお}きい\R{次数}{じ|すう}が2の\R{多項式}{た|こう|しき}を\textbf{2\R{次式}{じ|しき}}、1の\R{多項式}{た|こう|しき}を\textbf{1\R{次式}{じ|しき}}といいます。$3x^2-2x+5$ は2\R{次式}{じ|しき}、$2x+1$ は1\R{次式}{じ|しき}です（2-16の「1\R{次方程式}{じ|ほう|てい|しき}」の1\R{次}{じ}はこれ）。$h=20t-5t^2$ は $t$ の2\R{次式}{じ|しき}です。

\bkthrihead{\R{単項式}{たん|こう|しき}のかけ\R{算}{ざん}}
$3a^2b\times2ab$ は、×を\R{全部}{ぜん|ぶ}\R{書}{か}けば、ただのかけ\R{算}{ざん}のならびです。かけ\R{算}{ざん}は\R{順番}{じゅん|ばん}を\R{入}{い}れかえてよいので（\R{交換法則}{こう|かん|ほう|そく}、2-11）、\R{数}{かず}は\R{数}{かず}どうし、\R{同}{おな}じ\R{文字}{も|じ}は\R{同}{おな}じ\R{文字}{も|じ}どうしで\R{集}{あつ}めます。
\begin{align*}
3a^2b\times2ab&=3\times a\times a\times b\times2\times a\times b\\
&=(3\times2)\times(a\times a\times a)\times(b\times b)=6a^3b^2
\end{align*}
$a$ は $2+1=3$ こ、$b$ は $1+1=2$ こです。$a^2\times a=a^{2+1}=a^3$（2-12の\R{指数}{し|すう}のきまり）と\R{同}{おな}じです。\R{次数}{じ|すう}は、3\R{次}{じ}×2\R{次}{じ}で $3+2=5$ \R{次}{じ}になります。かけ\R{算}{ざん}すると、\R{次数}{じ|すう}はたし\R{算}{ざん}になるのです。

\bksubsection{2}{かっこ×かっこ：\R{全部}{ぜん|ぶ}の\R{組}{く}み\R{合}{あ}わせをかけてたす}

\leavevmode\label{bk:3-8:kumiawase}2-12では、$3(x+2)=3x+6$ のように、かっこの\R{外}{そと}の\R{数}{かず}を\R{中}{なか}のどちらにもかけました（\R{分配法則}{ぶん|ぱい|ほう|そく}）。では、かっこどうしのかけ\R{算}{ざん} $(a+b)(c+d)$ は、どうすればよいでしょう。2\R{通}{とお}りの\R{方法}{ほう|ほう}で\R{考}{かんが}えます。

\bkthrihead{\R{考}{かんが}え\R{方}{かた}A：\R{長方形}{ちょう|ほう|けい}の\R{面積}{めん|せき}で}
\noindent\begin{minipage}[c]{0.46\linewidth}\centering
\begin{tikzpicture}[x=1mm, y=1mm]
  \fill[bkBlue!14] (0,14) rectangle (32,36);
  \fill[bkOrange!25] (32,14) rectangle (48,36);
  \fill[bkGreen!22] (0,0) rectangle (32,14);
  \fill[bkPurple!22] (32,0) rectangle (48,14);
  \draw[bkInk, line width=0.8pt] (0,0) rectangle (48,36);
  \draw[bkInk, line width=0.6pt] (32,0) -- (32,36) (0,14) -- (48,14);
  \node[font=\small] at (16,25) {$ac$};
  \node[font=\small] at (40,25) {$bc$};
  \node[font=\small] at (16,7) {$ad$};
  \node[font=\small] at (40,7) {$bd$};
  \bkbrace{(0,37.5)}{(32,37.5)}{$a$}{1}\bkbrace{(32,37.5)}{(48,37.5)}{$b$}{1}
  \node[font=\small, anchor=east] at (-1.5,25) {$c$};
  \node[font=\small, anchor=east] at (-1.5,7) {$d$};
\end{tikzpicture}
\bkfigcaption{\R{横}{よこ} $a+b$、たて $c+d$ の\R{長方形}{ちょう|ほう|けい}}\label{bk:3-8:fig-cho}
\end{minipage}\hfill
\begin{minipage}[c]{0.5\linewidth}
\R{横}{よこ}が $a+b$、たてが $c+d$ の\R{長方形}{ちょう|ほう|けい}を、\R{横}{よこ}とたての\R{線}{せん}で4つに\R{分}{わ}けます。\R{全体}{ぜん|たい}の\R{面積}{めん|せき}は $(a+b)(c+d)$ です。4つの\R{小}{ちい}さな\R{長方形}{ちょう|ほう|けい}の\R{面積}{めん|せき}は $ac$、$bc$、$ad$、$bd$ です。\R{全体}{ぜん|たい}は4つの\R{合計}{ごう|けい}なので、
\[ (a+b)(c+d)=ac+ad+bc+bd \]
\end{minipage}

\bkthrihead{\R{考}{かんが}え\R{方}{かた}B：\R{分配法則}{ぶん|ぱい|ほう|そく}を2\R{回}{かい}\R{使}{つか}う}
$c+d$ を1つのかたまりと\R{見}{み}て、$M$ と\R{書}{か}きます。すると $(a+b)M$ で、2-11の\R{分配法則}{ぶん|ぱい|ほう|そく}がそのまま\R{使}{つか}えます。
\begin{enumerate}[label={\bkcircled[bkOrange]{\arabic*}}, leftmargin=2.4\zw, itemsep=2pt, topsep=3pt]
\item $M$ を $a$ にも $b$ にもかける：\quad $(a+b)M=aM+bM$
\item $M$ を $c+d$ にもどす：\quad $aM+bM=a(c+d)+b(c+d)$
\item もう1\R{回}{かい}、\R{分配法則}{ぶん|ぱい|ほう|そく}：\quad $a(c+d)+b(c+d)=ac+ad+bc+bd$
\end{enumerate}
\R{面積}{めん|せき}で\R{考}{かんが}えても、\R{分配法則}{ぶん|ぱい|ほう|そく}で\R{考}{かんが}えても、\R{同}{おな}じ4つの\R{項}{こう}になりました。\R{左}{ひだり}のかっこの\R{項}{こう}（$a$、$b$）の1つずつと、\R{右}{みぎ}のかっこの\R{項}{こう}（$c$、$d$）の1つずつを、\textbf{\R{全部}{ぜん|ぶ}の\R{組}{く}み\R{合}{あ}わせ}でかけて、たしています。$2\times2=4$ \R{組}{くみ}です。このように、かっこのある\R{式}{しき}をかけ\R{算}{ざん}して、かっこのない\R{多項式}{た|こう|しき}にすることを\textbf{\R{展開}{てん|かい}}といいます。\bkidx{てんかい}{\R{展開}{てん|かい}}

\begin{taisetsu}
かっこをかけるときは、\R{全部}{ぜん|ぶ}の\R{組}{く}み\R{合}{あ}わせをかけてたす。
\end{taisetsu}

\bkthrihead{\R{組}{く}み\R{合}{あ}わせの\R{表}{ひょう}を\R{使}{つか}う}
\R{項}{こう}が\R{多}{おお}いときや、マイナスの\R{項}{こう}があるときは、かけ\R{算}{ざん}の\R{表}{ひょう}を\R{作}{つく}ると、\R{組}{く}み\R{合}{あ}わせをわすれません（5-4でも\R{使}{つか}う\R{方法}{ほう|ほう}です）。$(2x+1)(x-3)$ なら、$x-3=x+(-3)$ と\R{見}{み}て、\R{表}{ひょう}を\R{作}{つく}ります。

\noindent\begin{minipage}[c]{0.46\linewidth}\centering
\begin{bktable}{colspec={Q[c,m,wd=16mm] Q[c,m,wd=16mm] Q[c,m,wd=16mm]}, rowsep=2pt, cell{1}{1}={bg=white}}
かける\R{数}{かず} & $\times\,x$ & $\times\,(-3)$ \\
$2x$ & $2x^2$ & $-6x$ \\
$1$ & $x$ & $-3$ \\
\end{bktable}
\end{minipage}\hfill
\begin{minipage}[c]{0.52\linewidth}
4つのますを\R{全部}{ぜん|ぶ}たして、\R{同類項}{どう|るい|こう}（$-6x$ と $x$）をまとめます（2-12）。
\[ 2x^2-6x+x-3=2x^2-5x-3 \]
\end{minipage}

\vspace{3pt}
\R{確}{たし}かめ：$x=2$ を\R{入}{い}れると、\R{左}{ひだり}は $(4+1)\times(2-3)=5\times(-1)=-5$、\R{右}{みぎ}は $8-10-3=-5$ で\R{合}{あ}います。

\begin{column}{かけ\R{算}{ざん}の\R{筆算}{ひっ|さん}も、\R{展開}{てん|かい}だった}
$23\times14$ の\R{筆算}{ひっ|さん}を、\R{展開}{てん|かい}で\R{見}{み}てみましょう。$23=20+3$、$14=10+4$ と\R{分}{わ}けると、
\[ (20+3)(10+4)=200+80+30+12=322 \]
\R{筆算}{ひっ|さん}の1\R{段}{だん}めの92は $23\times4=80+12$、2\R{段}{だん}めの230は $23\times10=200+30$ です。\R{小学校}{しょう|がっ|こう}から\R{使}{つか}ってきた\R{筆算}{ひっ|さん}は、「\R{全部}{ぜん|ぶ}の\R{組}{く}み\R{合}{あ}わせをかけてたす」を、\R{位}{くらい}ごとに\R{整理}{せい|り}して\R{書}{か}く\R{方法}{ほう|ほう}だったのです。
\end{column}

\bksubsection{3}{2\R{乗}{じょう}の\R{展開}{てん|かい}：$(a+b)^2$ と $(a-b)^2$}

\label{bk:3-8:nijou}\bkthrihead{$(a+b)^2$：やってみようの\R{正方形}{せい|ほう|けい}}
\noindent\begin{minipage}[c]{0.42\linewidth}\centering
\begin{tikzpicture}[x=1mm, y=1mm]
  \fill[bkBlue!14] (0,12) rectangle (34,46);
  \fill[bkOrange!25] (34,12) rectangle (46,46);
  \fill[bkOrange!25] (0,0) rectangle (34,12);
  \fill[bkRed!35] (34,0) rectangle (46,12);
  \draw[bkInk, line width=0.8pt] (0,0) rectangle (46,46);
  \draw[bkInk, line width=0.6pt] (34,0) -- (34,46) (0,12) -- (46,12);
  \node[font=\small] at (17,29) {$a^2$};
  \node[font=\small] at (40,29) {$ab$};
  \node[font=\small] at (17,6) {$ab$};
  \node[font=\small] at (40,6) {$b^2$};
  \bkbrace{(0,47.5)}{(34,47.5)}{$a$}{1}\bkbrace{(34,47.5)}{(46,47.5)}{$b$}{1}
  \node[font=\small, anchor=east] at (-1.5,29) {$a$};
  \node[font=\small, anchor=east] at (-1.5,6) {$b$};
\end{tikzpicture}
\bkfigcaption{1\R{辺}{ぺん} $a+b$ の\R{正方形}{せい|ほう|けい}}\label{bk:3-8:fig-sei}
\end{minipage}\hfill
\begin{minipage}[c]{0.55\linewidth}
やってみようの\R{正方形}{せい|ほう|けい}で、10を $a$、3を $b$ にかえます。1\R{辺}{ぺん} $a+b$ の\R{正方形}{せい|ほう|けい}は、\R{大}{おお}きな\R{正方形}{せい|ほう|けい} $a^2$（\R{青}{あお}）、\R{細長}{ほそ|なが}い\R{長方形}{ちょう|ほう|けい} $ab$ が2まい（オレンジ）、\R{小}{ちい}さな\R{正方形}{せい|ほう|けい} $b^2$（\R{赤}{あか}）に\R{分}{わ}かれます。
\[ (a+b)^2=a^2+2ab+b^2 \]
\end{minipage}

\vspace{3pt}
\R{式}{しき}でも\R{確}{たし}かめます。$(a+b)^2=(a+b)(a+b)$ なので、\R{組}{く}み\R{合}{あ}わせは $a\times a$、$a\times b$、$b\times a$、$b\times b$ の4つです。$ab$ と $ba$ は\R{同}{おな}じ（\R{交換法則}{こう|かん|ほう|そく}）なので、まとめて $2ab$ になります。まん\R{中}{なか}の「2」は、オレンジの\R{長方形}{ちょう|ほう|けい}が2まいあることです。

\bkthrihead{$(1+h)^2$：\R{第}{だい}5\R{巻}{かん}で\R{何度}{なん|ど}も\R{使}{つか}う\R{形}{かたち}}
$a=1$、$b=h$ とすると、
\bkthriwaku{$(1+h)^2=1+2h+h^2$}
\R{数}{かず}で\R{確}{たし}かめましょう。$h=0.1$ なら、\R{左}{ひだり}は $1.1^2=1.1\times1.1=1.21$、\R{右}{みぎ}は $1+0.2+0.01=1.21$ で、ぴったり\R{合}{あ}います。$h$ が\R{小}{ちい}さいと、$h^2$（\R{赤}{あか}い\R{小}{ちい}さな\R{正方形}{せい|ほう|けい}）はとても\R{小}{ちい}さくなります。$h=0.1$ なら0.01、$h=0.01$ なら0.0001です。この「$h^2$ はとても\R{小}{ちい}さい」ことを、5-2・5-8で\R{使}{つか}います。

\bkthrihead{$(a-b)^2$：\R{正方形}{せい|ほう|けい}から\R{切}{き}りとる}
\noindent\begin{minipage}[c]{0.42\linewidth}\centering
\begin{tikzpicture}[x=1mm, y=1mm]
  \fill[bkBlue!14] (0,0) rectangle (34,34);
  \fill[bkOrange!25] (34,0) rectangle (46,34);
  \fill[bkOrange!25] (0,34) rectangle (34,46);
  \fill[bkRed!35] (34,34) rectangle (46,46);
  \draw[bkInk, line width=0.8pt] (0,0) rectangle (46,46);
  \draw[bkInk, line width=0.6pt, dashed] (34,0) -- (34,46) (0,34) -- (46,34);
  \node[font=\small] at (17,17) {$(a-b)^2$};
  \node[font=\small, text=bkOrange!80!black] at (40,17) {\rotatebox{90}{$b(a-b)$}};
  \node[font=\small, text=bkOrange!80!black] at (17,40) {$b(a-b)$};
  \node[font=\small, text=bkRed!80!black] at (40,40) {$b^2$};
  \bkbrace{(0,-1.5)}{(34,-1.5)}{$a-b$}{-1}\bkbrace{(34,-1.5)}{(46,-1.5)}{$b$}{-1}
  \bkbrace{(0,47.5)}{(46,47.5)}{$a$}{1}
\end{tikzpicture}
\bkfigcaption{1\R{辺}{ぺん} $a$ の\R{正方形}{せい|ほう|けい}から、はばが $b$ の\R{帯}{おび}を\R{切}{き}りとる}\label{bk:3-8:fig-hiku}
\end{minipage}\hfill
\begin{minipage}[c]{0.55\linewidth}
1\R{辺}{ぺん} $a$ の\R{正方形}{せい|ほう|けい}（\R{面積}{めん|せき} $a^2$）から、\R{右}{みぎ}と\R{上}{うえ}の、はばが $b$ の\R{帯}{おび}を\R{切}{き}りとると、のこりの\R{青}{あお}い\R{正方形}{せい|ほう|けい}は1\R{辺}{ぺん} $a-b$ です。

\R{帯}{おび}は、1\R{本}{ぽん}が\R{長}{なが}さ $a$・はば $b$ で\R{面積}{めん|せき} $ab$。2\R{本}{ほん}ひくと $a^2-2ab$ です。ところが、\R{右上}{みぎ|うえ}の\R{赤}{あか}い\R{角}{かど} $b^2$ は、2\R{本}{ほん}の\R{帯}{おび}の\R{両方}{りょう|ほう}に\R{入}{はい}っているので、2\R{回}{かい}ひいたことになります。1\R{回}{かい}\R{分}{ぶん}を\R{返}{かえ}して、$+b^2$ します。
\end{minipage}
\bkthriline{$(a-b)^2=a^2-2ab+b^2$}
\R{式}{しき}なら、$(a-b)^2=\{a+(-b)\}^2$ と\R{見}{み}て $(a+b)^2$ の $b$ を $-b$ にかえれば\R{同}{おな}じです。$2a\times(-b)=-2ab$、$(-b)^2=b^2$（マイナス×マイナスはプラス、2-2）だからです。

\R{数}{かず}で\R{確}{たし}かめます。$a=10$、$b=3$ なら、\R{左}{ひだり}は $7^2=49$、\R{右}{みぎ}は $100-60+9=49$ で\R{合}{あ}います。

\bkthrihead{$(12-2x)^2$：\R{数}{かず}がついていても\R{同}{おな}じ}
$a=12$、$b=2x$ と\R{見}{み}ます。$b$ が $2x$ という\R{単項式}{たん|こう|しき}でも、かたまりとして\R{同}{おな}じように\R{使}{つか}えます。
\begin{align*}
(12-2x)^2&=12^2-2\times12\times2x+(2x)^2\\
&=144-48x+4x^2
\end{align*}
$(2x)^2=2x\times2x=4x^2$ です。$2x^2$ としないように\R{気}{き}をつけます（\R{道具}{どう|ぐ}1の\R{単項式}{たん|こう|しき}のかけ\R{算}{ざん}）。この\R{式}{しき}は、5-7で\R{箱}{はこ}の\R{体積}{たい|せき}がいちばん\R{大}{おお}きくなる\R{形}{かたち}を\R{求}{もと}めるときに\R{使}{つか}います。

\bksubsection{4}{\R{和}{わ}と\R{差}{さ}の\R{積}{せき}・$(x+a)(x+b)$}

\label{bk:3-8:wasa}\bkthrihead{$(a+b)(a-b)$：まん\R{中}{なか}が\R{消}{き}える}
\R{組}{く}み\R{合}{あ}わせの4つをかけてたします。
\[ (a+b)(a-b)=a^2-ab+ab-b^2=a^2-b^2 \]
$-ab$ と $+ab$ が\R{打}{う}ち\R{消}{け}し\R{合}{あ}って、2つの\R{項}{こう}だけがのこります。\R{図}{ず}でも\R{見}{み}てみましょう。

\begin{bkfigure}
\begin{tikzpicture}[x=0.85mm, y=0.85mm]
  % 左：1辺 a の正方形から、右上の角の b×b をとった形  % noruby
  \fill[bkBlue!14] (0,0) rectangle (40,30);
  \fill[bkOrange!25] (0,30) rectangle (30,40);
  \draw[bkRed!80!black, line width=0.6pt, dashed] (30,30) rectangle (40,40);
  \node[font=\small, text=bkRed!80!black] at (35,35) {$b^2$};
  \draw[bkInk, line width=0.8pt] (0,0) -- (40,0) -- (40,30) -- (30,30) -- (30,40) -- (0,40) -- cycle;
  \draw[bkInk, line width=0.5pt, dashed] (0,30) -- (30,30);
  \node[font=\small] at (20,15) {$a(a-b)$};
  \node[font=\small] at (15,35) {$b(a-b)$};
  \bkbrace{(0,-1.5)}{(40,-1.5)}{$a$}{-1}
  \node[font=\small, anchor=east] at (-1.5,20) {$a$};
  \node[font=\sffamily\footnotesize, anchor=north] at (20,-9) {\R{面積}{めん|せき} $a^2-b^2$};
  \draw[bk arrow, bkGray] (48,20) -- (60,20);
  \node[font=\sffamily\scriptsize, anchor=south, align=center] at (54,21.5) {オレンジを\\\R{回}{まわ}して\\\R{右}{みぎ}へ\R{移}{うつ}す};
  % 右：たて a-b、横 a+b の長方形  % noruby
  \begin{scope}[shift={(66,10)}]
  \fill[bkBlue!14] (0,0) rectangle (40,30);
  \fill[bkOrange!25] (40,0) rectangle (50,30);
  \draw[bkInk, line width=0.8pt] (0,0) rectangle (50,30);
  \draw[bkInk, line width=0.5pt, dashed] (40,0) -- (40,30);
  \node[font=\small] at (20,15) {$a(a-b)$};
  \node[font=\small] at (45,15) {\rotatebox{90}{$b(a-b)$}};
  \bkbrace{(0,-1.5)}{(40,-1.5)}{$a$}{-1}\bkbrace{(40,-1.5)}{(50,-1.5)}{$b$}{-1}
  \node[font=\small, anchor=west] at (51.5,15) {$a-b$};
  \node[font=\sffamily\footnotesize, anchor=north] at (25,-9) {$(a+b)(a-b)$};
  \end{scope}
\end{tikzpicture}
\bkfigcaption{1\R{辺}{ぺん} $a$ の\R{正方形}{せい|ほう|けい}から\R{角}{かど}の $b^2$ をとり、オレンジの\R{部分}{ぶ|ぶん}を\R{回}{まわ}して\R{右}{みぎ}へ\R{移}{うつ}すと、\R{横}{よこ} $a+b$・たて $a-b$ の\R{長方形}{ちょう|ほう|けい}}\label{bk:3-8:fig-wasa}
\end{bkfigure}

\R{左}{ひだり}の\R{形}{かたち}は、1\R{辺}{ぺん} $a$ の\R{正方形}{せい|ほう|けい}から\R{右上}{みぎ|うえ}の $b^2$ をとったので、\R{面積}{めん|せき}は $a^2-b^2$ です。\R{上}{うえ}のオレンジの\R{長方形}{ちょう|ほう|けい}（\R{横}{よこ} $a-b$、たて $b$）を90°\R{回}{まわ}して\R{立}{た}て、\R{青}{あお}い\R{長方形}{ちょう|ほう|けい}（\R{横}{よこ} $a$、たて $a-b$）の\R{右}{みぎ}にならべると、\R{横}{よこ} $a+b$・たて $a-b$ の\R{長方形}{ちょう|ほう|けい}になります。\R{面積}{めん|せき}は\R{変}{か}わらないので、$(a+b)(a-b)=a^2-b^2$ です。

\R{暗算}{あん|ざん}にも\R{使}{つか}えます。$101\times99=(100+1)(100-1)=10000-1=9999$。

\bkthrihead{$(x+a)(x+b)$：\R{和}{わ}と\R{積}{せき}}
\noindent\begin{minipage}[c]{0.42\linewidth}\centering
\begin{tikzpicture}[x=1mm, y=1mm]
  \fill[bkBlue!14] (0,10) rectangle (30,40);
  \fill[bkOrange!25] (30,10) rectangle (38,40);
  \fill[bkGreen!22] (0,0) rectangle (30,10);
  \fill[bkPurple!22] (30,0) rectangle (38,10);
  \draw[bkInk, line width=0.8pt] (0,0) rectangle (38,40);
  \draw[bkInk, line width=0.6pt] (30,0) -- (30,40) (0,10) -- (38,10);
  \node[font=\small] at (15,25) {$x^2$};
  \node[font=\small] at (34,25) {$ax$};
  \node[font=\small] at (15,5) {$bx$};
  \node[font=\small] at (34,5) {$ab$};
  \bkbrace{(0,41.5)}{(30,41.5)}{$x$}{1}\bkbrace{(30,41.5)}{(38,41.5)}{$a$}{1}
  \node[font=\small, anchor=east] at (-1.5,25) {$x$};
  \node[font=\small, anchor=east] at (-1.5,5) {$b$};
\end{tikzpicture}
\end{minipage}\hfill
\begin{minipage}[c]{0.55\linewidth}
\R{横}{よこ} $x+a$、たて $x+b$ の\R{長方形}{ちょう|ほう|けい}です。$ax$ と $bx$ は\R{同類項}{どう|るい|こう}なので、まとめて $(a+b)x$ です（2-12）。
\[ (x+a)(x+b)=x^2+(a+b)x+ab \]
$x$ の\R{係数}{けい|すう}は $a$ と $b$ の\textbf{\R{和}{わ}}、\R{数}{かず}だけの\R{項}{こう}は $a$ と $b$ の\textbf{\R{積}{せき}}です。
\end{minipage}

\vspace{3pt}
たとえば $(x+2)(x+5)=x^2+7x+10$ です（$2+5=7$、$2\times5=10$）。$a$ や $b$ がマイナスでも、\R{和}{わ}と\R{積}{せき}を\R{計算}{けい|さん}すれば\R{同}{おな}じです。$(t+5)(t-3)$ なら、$5+(-3)=2$、$5\times(-3)=-15$ なので、$t^2+2t-15$ です。

\begin{bkcenter}
\begin{bktable}{colspec={Q[l,m,wd=30mm] Q[l,m,wd=46mm] Q[l,m,wd=40mm]}, rowsep=1.5pt, cells={font=\small}, row{1}={font=\sffamily\bfseries\small}}
かっこの\R{形}{かたち} & \R{展開}{てん|かい}した\R{形}{かたち} & \R{図}{ず}で\R{見}{み}ると \\
$(a+b)^2$ & $a^2+2ab+b^2$ & \R{正方形}{せい|ほう|けい}を4つに\R{分}{わ}ける \\
$(a-b)^2$ & $a^2-2ab+b^2$ & \R{帯}{おび}を2\R{本}{ほん}\R{切}{き}り、\R{角}{かど}を\R{返}{かえ}す \\
$(a+b)(a-b)$ & $a^2-b^2$ & \R{角}{かど}をとって\R{並}{なら}べかえる \\
$(x+a)(x+b)$ & $x^2+(a+b)x+ab$ & \R{長方形}{ちょう|ほう|けい}を4つに\R{分}{わ}ける \\
\end{bktable}
\bktabcaption{よく\R{使}{つか}う\R{展開}{てん|かい}の\R{形}{かたち}。どれも「\R{全部}{ぜん|ぶ}の\R{組}{く}み\R{合}{あ}わせをかけてたす」から\R{作}{つく}れる。}
\end{bkcenter}

\R{表}{ひょう}の\R{形}{かたち}は、\R{覚}{おぼ}えていなくてもかまいません。わすれたら、\R{組}{く}み\R{合}{あ}わせの\R{表}{ひょう}や\R{面積}{めん|せき}の\R{図}{ず}から、いつでも\R{作}{つく}り\R{直}{なお}せます。

\bksubsection{5}{\R{立方体}{りっ|ぽう|たい}で\R{見}{み}る $(a+b)^3$}

\leavevmode\label{bk:3-8:cube}3\R{乗}{じょう}も、2\R{乗}{じょう}を\R{使}{つか}って\R{展開}{てん|かい}できます。
\[ (a+b)^3=(a+b)^2\times(a+b)=(a^2+2ab+b^2)(a+b) \]
\R{左}{ひだり}のかっこの3つの\R{項}{こう}と、\R{右}{みぎ}のかっこの2つの\R{項}{こう}で、$3\times2=6$ \R{組}{くみ}の\R{組}{く}み\R{合}{あ}わせができます。

\begin{bkcenter}
\begin{bktable}{colspec={Q[c,m,wd=24mm] Q[c,m,wd=24mm] Q[c,m,wd=24mm]}, rowsep=2pt, cell{1}{1}={bg=white}}
かける\R{数}{かず} & $\times\,a$ & $\times\,b$ \\
$a^2$ & $a^3$ & $a^2b$ \\
$2ab$ & $2a^2b$ & $2ab^2$ \\
$b^2$ & $ab^2$ & $b^3$ \\
\end{bktable}
\bktabcaption{6つの\R{組}{く}み\R{合}{あ}わせ。$a^2b$ は $1+2=3$ こ、$ab^2$ は $2+1=3$ こある。}
\end{bkcenter}
\bkthriline{$(a+b)^3=a^3+3a^2b+3ab^2+b^3$}
\R{数}{かず}で\R{確}{たし}かめます。$a=10$、$b=1$ なら、\R{左}{ひだり}は $11^3=1331$ です。\R{右}{みぎ}は $1000+300+30+1$ で、やはり1331です。

$(a+b)^3$ は、1\R{辺}{ぺん} $a+b$ の\R{立方体}{りっ|ぽう|たい}の\R{体積}{たい|せき}です。1\R{辺}{ぺん} $a$ の\R{立方体}{りっ|ぽう|たい}（\R{青}{あお}）のたて・\R{横}{よこ}・\R{高}{たか}さを $b$ ずつのばし、ふえた\R{部分}{ぶ|ぶん}を\R{引}{ひ}きはなして、\R{種類}{しゅ|るい}ごとに\R{並}{なら}べます。

\begin{bkfigure}
\begin{tikzpicture}[x={(1mm,0mm)}, y={(0.48mm,0.32mm)}, z={(0mm,1mm)}]
  \begin{scope}
    \bkthricube{16}{6}{0}{0}{0}
    \node[font=\sffamily\small, anchor=north] at (8,-6,-2.2) {$a$};
    \node[font=\sffamily\small, anchor=north] at (19,-6,-2.2) {$b$};
  \end{scope}
  \draw[bk arrow, bkGray] (32,0,12) -- (44,0,12);
  \node[font=\sffamily\scriptsize, anchor=south, align=center] at (38,0,13.5) {\R{引}{ひ}き\\はなす};
  \begin{scope}[shift={(54,0,0)}]
    \bkthricube{16}{6}{8}{18}{8}
  \end{scope}
\end{tikzpicture}
\bkfigcaption{\R{左}{ひだり}：1\R{辺}{ぺん} $a+b$ の\R{立方体}{りっ|ぽう|たい}。\R{右}{みぎ}：\R{色}{いろ}ごとに\R{引}{ひ}きはなした\R{形}{かたち}（5-4の\R{図}{ず}と\R{同}{おな}じ\R{色}{いろ}\R{分}{わ}け）}\label{bk:3-8:fig-cube1}
\end{bkfigure}

\begin{bkfigure}
\begin{tikzpicture}[x={(1mm,0mm)}, y={(0.48mm,0.32mm)}, z={(0mm,1mm)}, lb/.style={font=\sffamily\footnotesize, anchor=north, align=center}]
  \bkthribox{0}{0}{0}{16}{16}{16}{bkBlue}
  \foreach \k in {2,1,0} {\bkthribox{30}{0}{\k*9}{16}{16}{5}{bkOrange}}
  \foreach \k in {2,1,0} {\bkthribox{62}{0}{\k*9}{16}{5}{5}{bkGreen}}
  \bkthribox{94}{0}{0}{5}{5}{5}{bkPurple}
  \node[lb] at (10,0,-2.5) {\textbf{\R{立方体}{りっ|ぽう|たい} 1こ}\\$a^3$};
  \node[lb] at (42,0,-2.5) {\textbf{\R{板}{いた} 3まい}\\1まい $a\times a\times b$\\\R{合計}{ごう|けい} $3a^2b$};
  \node[lb] at (71,0,-2.5) {\textbf{\R{角柱}{かく|ちゅう} 3\R{本}{ぼん}}\\1\R{本}{ぽん} $a\times b\times b$\\\R{合計}{ごう|けい} $3ab^2$};
  \node[lb] at (98,0,-2.5) {\textbf{\R{小}{ちい}さな}\\\textbf{\R{立方体}{りっ|ぽう|たい} 1こ}\\$b^3$};
\end{tikzpicture}
\bkfigcaption{\R{種類}{しゅ|るい}ごとに\R{並}{なら}べた。\R{板}{いた}が3まい、\R{角柱}{かく|ちゅう}が3\R{本}{ぼん}あるのが、\R{式}{しき}の「3」}\label{bk:3-8:fig-cube2}
\end{bkfigure}

\R{青}{あお}い\R{立方体}{りっ|ぽう|たい} $a^3$、うすい\R{板}{いた}（オレンジ）が3まいで $3a^2b$、\R{細}{ほそ}い\R{角柱}{かく|ちゅう}（\R{緑}{みどり}）が3\R{本}{ぼん}で $3ab^2$、\R{小}{ちい}さな\R{立方体}{りっ|ぽう|たい}（\R{紫}{むらさき}）が $b^3$ です。\R{表}{ひょう}で\R{求}{もと}めた\R{式}{しき}と、ぴったり\R{同}{おな}じです。5-4では、この\R{図}{ず}を\R{使}{つか}って、$x^3$ の\R{微分}{び|ぶん}のきまりを\R{確}{たし}かめます。

\bksubsection{6}{\R{共通因数}{きょう|つう|いん|すう}でくくる：\R{展開}{てん|かい}の\R{逆}{ぎゃく}}

\leavevmode\label{bk:3-8:kukuru}\R{展開}{てん|かい}は、かっこをはずす\R{計算}{けい|さん}でした。こんどは\R{反対}{はん|たい}に、かっこを\R{作}{つく}ります。

\R{分配法則}{ぶん|ぱい|ほう|そく} $c(a+b)=ca+cb$ を、\R{右}{みぎ}から\R{左}{ひだり}へ\R{読}{よ}みます。$ca+cb$ のどちらの\R{項}{こう}にも $c$ がかかっているので、$c$ を\R{外}{そと}に\R{出}{だ}して $c(a+b)$ とまとめられます。どの\R{項}{こう}にもかかっている\R{数}{かず}や\R{文字}{も|じ}を\textbf{\R{共通因数}{きょう|つう|いん|すう}}といい、それを\R{外}{そと}に\R{出}{だ}すことを「\textbf{くくる}」といいます。

\noindent\begin{minipage}[c]{0.4\linewidth}\centering
\begin{tikzpicture}[x=1mm, y=1mm]
  \fill[bkBlue!14] (0,0) rectangle (24,14);
  \fill[bkOrange!25] (24,0) rectangle (40,14);
  \draw[bkInk, line width=0.8pt] (0,0) rectangle (40,14);
  \draw[bkInk, line width=0.6pt] (24,0) -- (24,14);
  \node[font=\small] at (12,7) {$ca$};
  \node[font=\small] at (32,7) {$cb$};
  \bkbrace{(0,15.5)}{(24,15.5)}{$a$}{1}\bkbrace{(24,15.5)}{(40,15.5)}{$b$}{1}
  \node[font=\small, anchor=east, text=bkRed!80!black] at (-1.5,7) {$c$};
\end{tikzpicture}
\end{minipage}\hfill
\begin{minipage}[c]{0.57\linewidth}
\R{面積}{めん|せき} $ca$ と $cb$ の2まいの\R{長方形}{ちょう|ほう|けい}は、たての\R{長}{なが}さ $c$ が\R{同}{おな}じです。だから\R{横}{よこ}にならべてくっつけると、たて $c$・\R{横}{よこ} $a+b$ の1まいの\R{長方形}{ちょう|ほう|けい}になります。\R{同}{おな}じたての\R{長}{なが}さが、\R{共通因数}{きょう|つう|いん|すう}です。
\end{minipage}

\bkthrihead{12でくくる}
$12x^2-96x+144$ は、どの\R{項}{こう}も12でわりきれます。$96=12\times8$、$144=12\times12$ だからです。\R{項}{こう}ごとに「12×\R{何}{なに}か」の\R{形}{かたち}にしてから、12を\R{外}{そと}に\R{出}{だ}します。
\begin{align*}
12x^2-96x+144&=12\times x^2-12\times8x+12\times12\\
&=12(x^2-8x+12)
\end{align*}
\R{確}{たし}かめは\R{展開}{てん|かい}です。$12(x^2-8x+12)$ のかっこをはずすと、もとの\R{式}{しき}にもどります。5-7では、このように12でくくってから\R{計算}{けい|さん}を\R{進}{すす}めます。

\bkthrihead{\R{文字}{も|じ}でもくくれる}
$2x^2-8x$ は、$2x^2=2x\times x$、$8x=2x\times4$ なので、どちらの\R{項}{こう}にも $2x$ がかかっています。$2x^2-8x=2x(x-4)$ です。\R{数}{かず}だけの2でくくって $2(x^2-4x)$ としてもまちがいではありませんが、かっこの\R{中}{なか}にまだ $x$ がのこっています。\R{共通因数}{きょう|つう|いん|すう}は、できるだけ\R{全部}{ぜん|ぶ}くくり\R{出}{だ}します。

\bkthrihead{\R{展開}{てん|かい}を\R{逆}{ぎゃく}に\R{読}{よ}むと：\R{因数分解}{いん|すう|ぶん|かい}の\R{入口}{いり|ぐち}}
\R{道具}{どう|ぐ}4の\R{展開}{てん|かい}
\[ (x+2)(x+5)=x^2+7x+10 \]
を、\R{右}{みぎ}から\R{左}{ひだり}へ\R{読}{よ}むと、$x^2+7x+10$ は $(x+2)(x+5)$ と\R{書}{か}けることになります。\R{和}{わ}が7、\R{積}{せき}が10の2つの\R{数}{かず}（2と5）が\R{見}{み}つかれば、かっこのかけ\R{算}{ざん}にもどせます。このように、\R{多項式}{た|こう|しき}をかけ\R{算}{ざん}の\R{形}{かたち}に\R{直}{なお}すことを\textbf{\R{因数分解}{いん|すう|ぶん|かい}}といいます。くわしくは3-10で、2\R{次方程式}{じ|ほう|てい|しき}を\R{解}{と}くときに\R{学}{まな}びます。

\bksubsection{7}{どんな $x$ でも\R{成}{な}り\R{立}{た}つ\R{等式}{とう|しき}：\R{恒等式}{こう|とう|しき}}

\leavevmode\label{bk:3-8:koutou}2-11では、$a+b=b+a$ や $(a+b)\times c=a\times c+b\times c$ のように、\R{文字}{も|じ}を\R{使}{つか}って「どんな\R{数}{かず}でも\R{成}{な}り\R{立}{た}つきまり」を1\R{行}{ぎょう}で\R{書}{か}きました。この\R{章}{しょう}で\R{作}{つく}った\R{展開}{てん|かい}の\R{式}{しき}も、じつは\R{同}{おな}じなかまです。ここでは、それを\R{方程式}{ほう|てい|しき}とくらべて、はっきりさせます。

\bkthristep{1}{\R{数}{かず}を\R{入}{い}れて、\R{左}{ひだり}と\R{右}{みぎ}をくらべる}
\R{展開}{てん|かい}した\R{式}{しき} $(x+1)^2=x^2+2x+1$ の $x$ に、いろいろな\R{数}{かず}を\R{入}{い}れて、「＝」の\R{左}{ひだり}（\R{左辺}{さ|へん}）と\R{右}{みぎ}（\R{右辺}{う|へん}）を\R{別々}{べつ|べつ}に\R{計算}{けい|さん}します（2-13）。

\begin{bkcenter}
\begin{bktable}{colspec={Q[c,m,wd=34mm] *{5}{Q[c,m,wd=11mm]}}, rowsep=1.5pt, cells={font=\small}}
$x$ & 0 & 1 & 2 & $-1$ & 0.5 \\
\R{左辺}{さ|へん} $(x+1)^2$ & 1 & 4 & 9 & 0 & 2.25 \\
\R{右辺}{う|へん} $x^2+2x+1$ & 1 & 4 & 9 & 0 & 2.25 \\
\SetCell{bg=bkYellow!30} \R{左右}{さ|ゆう}は？ & \R{同}{おな}じ & \R{同}{おな}じ & \R{同}{おな}じ & \R{同}{おな}じ & \R{同}{おな}じ \\
\end{bktable}
\bktabcaption{$(x+1)^2=x^2+2x+1$。どの $x$ を\R{入}{い}れても、\R{左右}{さ|ゆう}が\R{同}{おな}じ\R{値}{あたい}になる。}
\end{bkcenter}

たとえば $x=2$ なら、\R{左辺}{さ|へん}は $(2+1)^2=3^2=9$、\R{右辺}{う|へん}は $2^2+2\times2+1$ で、$4+4+1=9$ です。$x=0.5$ なら、\R{左辺}{さ|へん}は $1.5^2=2.25$、\R{右辺}{う|へん}は $0.25+1+1=2.25$ です。

\bkthristep{2}{\R{方程式}{ほう|てい|しき}だと、どうなる？}
こんどは $(x+1)^2=4$ で、\R{同}{おな}じことをします。

\begin{bkcenter}
\begin{bktable}{colspec={Q[c,m,wd=22mm] *{6}{Q[c,m,wd=12.5mm]}}, rowsep=1.5pt, cells={font=\small}}
$x$ & 0 & 1 & 2 & $-1$ & 0.5 & $-3$ \\
\R{左辺}{さ|へん} $(x+1)^2$ & 1 & 4 & 9 & 0 & 2.25 & 4 \\
\R{右辺}{う|へん} 4 & 4 & 4 & 4 & 4 & 4 & 4 \\
\SetCell{bg=bkYellow!30} \R{左右}{さ|ゆう}は？ & ちがう & \SetCell{bg=bkYellow!40}\textbf{\R{同}{おな}じ} & ちがう & ちがう & ちがう & \SetCell{bg=bkYellow!40}\textbf{\R{同}{おな}じ} \\
\end{bktable}
\bktabcaption{$(x+1)^2=4$。\R{左右}{さ|ゆう}が\R{同}{おな}じになるのは、$x=1$ と $x=-3$ のときだけ。}
\end{bkcenter}

\R{同}{おな}じになるのは、$x=1$ と $x=-3$ のときだけです。2\R{乗}{じょう}して4になる\R{数}{かず}は2と $-2$（2-8）なので、$x+1=2$ か $x+1=-2$、つまり $x=1$ か $x=-3$ で、ほかにはありません。

\bkthristep{3}{\R{名前}{な|まえ}をつける}
$(x+1)^2=x^2+2x+1$ のように、\textbf{\R{文字}{も|じ}にどんな\R{数}{かず}を\R{入}{い}れても\R{成}{な}り\R{立}{た}つ\R{等式}{とう|しき}}を、\textbf{\R{恒等式}{こう|とう|しき}}といいます。\bkidx{こうとうしき}{\R{恒等式}{こう|とう|しき}}「\R{恒}{こう}」は「いつも\R{変}{か}わらない」という\R{意味}{い|み}です。$(x+1)^2=4$ のように、\textbf{$x$ の\R{値}{あたい}によって、\R{成}{な}り\R{立}{た}ったり\R{成}{な}り\R{立}{た}たなかったりする\R{等式}{とう|しき}}は、2-16で\R{学}{まな}んだ\textbf{\R{方程式}{ほう|てい|しき}}です。$(x+1)^2=x^2+2x+5$ のように、\R{成}{な}り\R{立}{た}つ $x$ が1つもない\R{方程式}{ほう|てい|しき}もあります（\R{展開}{てん|かい}すると $1=5$ になってしまう。3-6の、\R{平行}{へい|こう}な2\R{直線}{ちょく|せん}の「\R{解}{かい}なし」と\R{同}{おな}じ）。

\bkthristep{4}{\R{天}{てん}びんで\R{見}{み}る}
\R{等式}{とう|しき}は\R{天}{てん}びんでした（2-14）。\R{左}{ひだり}の\R{皿}{さら}に\R{左辺}{さ|へん}の\R{値}{あたい}、\R{右}{みぎ}の\R{皿}{さら}に\R{右辺}{う|へん}の\R{値}{あたい}の\R{重}{おも}さをのせます。$x$ を\R{変}{か}えると、のせる\R{重}{おも}さが\R{変}{か}わります。

\begin{bkfigure}
\begin{tikzpicture}[x=0.5mm, y=0.5mm]
  \node[font=\sffamily\bfseries\small, anchor=west] at (-40,62) {\R{恒等式}{こう|とう|しき} $(x+1)^2=x^2+2x+1$};
  \foreach \cx/\xv/\v in {0/0/1, 80/1/4, 160/2/9} {
    \bkthribal{(\cx,0)}{0}{\v}{\v}
    \node[font=\sffamily\small, anchor=north] at (\cx,-3) {$x=\xv$};}
  \node[font=\sffamily\small, text=bkBlue!80!black, anchor=west, align=left] at (198,24) {いつも\\つり\R{合}{あ}う};
\end{tikzpicture}

\vspace{8pt}
\begin{tikzpicture}[x=0.5mm, y=0.5mm]
  \node[font=\sffamily\bfseries\small, anchor=west] at (-40,62) {\R{方程式}{ほう|てい|しき} $(x+1)^2=4$};
  \foreach \cx/\xv/\v/\ang in {0/0/1/-9, 80/1/4/0, 160/2/9/9} {
    \bkthribal{(\cx,0)}{\ang}{\v}{4}
    \node[font=\sffamily\small, anchor=north] at (\cx,-3) {$x=\xv$};}
  \node[font=\sffamily\small, text=bkRed!80!black, anchor=west, align=left] at (198,24) {$x=1$ と\\$x=-3$\\でだけ};
\end{tikzpicture}
\bkfigcaption{\R{左}{ひだり}の\R{皿}{さら}（\R{青}{あお}）が\R{左辺}{さ|へん}、\R{右}{みぎ}の\R{皿}{さら}（オレンジ）が\R{右辺}{う|へん}。\R{恒等式}{こう|とう|しき}は、$x$ を\R{変}{か}えてもずっとつり\R{合}{あ}う。}\label{bk:3-8:fig-tenbin}
\end{bkfigure}

\R{恒等式}{こう|とう|しき}の\R{天}{てん}びんは、$x$ をどんなに\R{変}{か}えても、いつもまっすぐです。\R{左右}{さ|ゆう}の\R{重}{おも}さが、いつもいっしょに\R{同}{おな}じだけ\R{変}{か}わるからです。\R{方程式}{ほう|てい|しき}の\R{天}{てん}びんは、$x=1$（と $x=-3$）のときだけつり\R{合}{あ}い、ほかの $x$ ではかたむきます。

\bkthristep{5}{グラフで\R{見}{み}る}
\R{左辺}{さ|へん}と\R{右辺}{う|へん}を、それぞれ $y$ として2\R{本}{ほん}のグラフにかくと、ちがいがもっとはっきりします。

\noindent\begin{minipage}[t]{0.49\linewidth}\centering\vspace{0pt}
\begin{tikzpicture}
\begin{axis}[bkplot, width=56mm, height=50mm, xmin=-4.3, xmax=2.6, ymin=-0.8, ymax=10.5, xtick={-3,-2,-1,1,2}, ytick={2,4,6,8}, xlabel={$x$}, ylabel={$y$}]
  \addplot[bkcurve, domain=-4.16:2.16, samples=100, line width=4pt, opacity=0.35] {(x+1)^2};
  \addplot[bkOrange, line width=1.2pt, dashed, domain=-4.16:2.16, samples=100] {x^2+2*x+1};
\end{axis}
\end{tikzpicture}
\bkfigcaption{$y=(x+1)^2$（\R{青}{あお}）と $y=x^2+2x+1$（\R{点線}{てん|せん}）は、ぴったり\R{重}{かさ}なる}
\end{minipage}\hfill
\begin{minipage}[t]{0.49\linewidth}\centering\vspace{0pt}
\begin{tikzpicture}
\begin{axis}[bkplot, width=56mm, height=50mm, xmin=-4.3, xmax=2.6, ymin=-0.8, ymax=10.5, xtick={-3,-2,-1,1,2}, ytick={2,4,6,8}, xlabel={$x$}, ylabel={$y$}]
  \addplot[bkcurve, domain=-4.16:2.16, samples=100] {(x+1)^2};
  \addplot[bkcurve2, domain=-4.2:2.5] {4};
  \addplot[only marks, mark=*, mark size=2pt, bkRed] coordinates {(-3,4) (1,4)};
\end{axis}
\end{tikzpicture}
\bkfigcaption{$y=(x+1)^2$ と $y=4$ が\R{交}{まじ}わるのは2\R{点}{てん}だけ}
\end{minipage}

\vspace{3pt}
\R{恒等式}{こう|とう|しき}では、2\R{本}{ほん}のグラフがどこでもぴったり\R{重}{かさ}なり、1\R{本}{ぽん}に\R{見}{み}えます。\R{左辺}{さ|へん}と\R{右辺}{う|へん}は、\R{見}{み}た\R{目}{め}がちがうだけで、\R{同}{おな}じ\R{関数}{かん|すう}（3-3）なのです。\R{方程式}{ほう|てい|しき}では、2\R{本}{ほん}のグラフは\R{交}{まじ}わる\R{点}{てん}（\R{赤}{あか}）でだけ\R{出}{で}あいます（3-6で\R{見}{み}た「\R{交点}{こう|てん}」と\R{同}{おな}じ）。

\begin{taisetsu}
\R{恒等式}{こう|とう|しき}：\R{文字}{も|じ}にどんな\R{数}{かず}を\R{入}{い}れても\R{成}{な}り\R{立}{た}つ\R{等式}{とう|しき}。\par
\hspace*{2\zw}（\R{例}{れい}）$(x+1)^2=x^2+2x+1$\par\vspace{2pt}
\R{方程式}{ほう|てい|しき}：$x$ の\R{値}{あたい}によって、\R{成}{な}り\R{立}{た}ったり\R{成}{な}り\R{立}{た}たなかったりする\R{等式}{とう|しき}。\R{成}{な}り\R{立}{た}つ $x$ がないこともある（3-6の\R{平行}{へい|こう}な2\R{直線}{ちょく|せん}）。\par
\hspace*{2\zw}（\R{例}{れい}）$(x+1)^2=4$
\end{taisetsu}

\R{恒等式}{こう|とう|しき}かどうかを\R{調}{しら}べるには、\R{数}{かず}を\R{入}{い}れて1つでも\R{合}{あ}わない $x$ があれば「\R{恒等式}{こう|とう|しき}ではない」とわかります。でも、いくつかの $x$ で\R{合}{あ}っただけでは、たまたまかもしれません。\R{恒等式}{こう|とう|しき}だと\R{言}{い}いきるには、\R{展開}{てん|かい}して\R{左右}{さ|ゆう}がまったく\R{同}{おな}じ\R{式}{しき}になることを\R{確}{たし}かめます。

\bksubsection{8}{\R{係数}{けい|すう}を\R{比}{くら}べる}

\leavevmode\label{bk:3-8:hikaku}\R{恒等式}{こう|とう|しき}を\R{使}{つか}うと、わからない\R{数}{かず}を\R{決}{き}めることができます。\R{次}{つぎ}の\R{問題}{もん|だい}を\R{考}{かんが}えましょう。

\bkthriwaku{$x^2+6x+10=(x+p)^2+q$ が\R{恒等式}{こう|とう|しき}になるような、\R{数}{かず} $p$ と $q$ は？}

この\R{形}{かたち}は、あとで\R{平方完成}{へい|ほう|かん|せい}（\R{道具}{どう|ぐ}10）に\R{使}{つか}います。

\bkthristep{1}{\R{右辺}{う|へん}を\R{展開}{てん|かい}する}
$(x+p)^2=x^2+2px+p^2$（\R{道具}{どう|ぐ}3の $(a+b)^2$ で、$a=x$、$b=p$）なので、
\[ (x+p)^2+q=x^2+2p\,x+(p^2+q) \]
$x$ にかかっている\R{数}{かず}は $2p$、\R{数}{かず}だけの\R{項}{こう}は $p^2+q$ です。

\bkthristep{2}{\R{同}{おな}じ\R{次数}{じ|すう}の\R{項}{こう}どうしを\R{比}{くら}べる}
\R{左辺}{さ|へん}と\R{右辺}{う|へん}を、\R{次数}{じ|すう}ごとに\R{上下}{じょう|げ}にならべます。

\begin{bkcenter}
\begin{bktable}{colspec={Q[c,m,wd=20mm] Q[c,m,wd=22mm] Q[c,m,wd=22mm] Q[c,m,wd=26mm]}, rowsep=1.5pt, cells={font=\small}}
 & $x^2$ の\R{係数}{けい|すう} & $x$ の\R{係数}{けい|すう} & \R{数}{かず}だけの\R{項}{こう} \\
\R{左辺}{さ|へん} & 1 & 6 & 10 \\
\R{右辺}{う|へん} & 1 & $2p$ & $p^2+q$ \\
\end{bktable}
\bktabcaption{\R{次数}{じ|すう}ごとに\R{係数}{けい|すう}を\R{比}{くら}べる}
\end{bkcenter}

\R{恒等式}{こう|とう|しき}になるのは、\R{左右}{さ|ゆう}がまったく\R{同}{おな}じ\R{式}{しき}になるとき、つまり\R{同}{おな}じ\R{次数}{じ|すう}どうしの\R{係数}{けい|すう}がそれぞれ\R{等}{ひと}しいときです。
\begin{itemize}[label={\color{bkOrange}●}, leftmargin=1.6\zw, itemsep=2pt, topsep=3pt]
\item $x$ の\R{係数}{けい|すう}を\R{比}{くら}べて $2p=6$。だから $p=3$。
\item \R{数}{かず}だけの\R{項}{こう}を\R{比}{くら}べて $p^2+q=10$。ここに $p=3$ を\R{入}{い}れて $9+q=10$。だから $q=1$。
\end{itemize}
だから
\[ x^2+6x+10=(x+3)^2+1 \]
です。\R{右辺}{う|へん}を\R{展開}{てん|かい}すると $x^2+6x+9+1$ で、たしかに\R{左辺}{さ|へん}と\R{同}{おな}じです。このように、\R{恒等式}{こう|とう|しき}の\R{両辺}{りょう|へん}で\R{同}{おな}じ\R{次数}{じ|すう}の\R{係数}{けい|すう}を\R{等}{ひと}しいとおいて、わからない\R{数}{かず}を\R{求}{もと}める\R{方法}{ほう|ほう}を、\textbf{\R{係数}{けい|すう}を\R{比}{くら}べる}（\R{係数}{けい|すう}\R{比較}{ひ|かく}）といいます。\bkidx{けいすうをくらべる}{\R{係数}{けい|すう}を\R{比}{くら}べる}

\bkthristep{3}{なぜ\R{比}{くら}べてよいの？}
「\R{係数}{けい|すう}が\R{等}{ひと}しければ\R{同}{おな}じ\R{式}{しき}」は\R{当}{あ}たり\R{前}{まえ}です。では\R{反対}{はん|たい}に、「どんな $x$ でも\R{等}{ひと}しいなら、\R{係数}{けい|すう}も\R{等}{ひと}しい」と\R{言}{い}えるでしょうか。1\R{次式}{じ|しき}で\R{確}{たし}かめます。

$ax+b=a'x+b'$ が、どんな $x$ でも\R{成}{な}り\R{立}{た}つとします（$a'$ は「エー・ダッシュ」と\R{読}{よ}み、$a$ とは\R{別}{べつ}の\R{数}{かず}です）。\R{右辺}{う|へん}を\R{左}{ひだり}へ\R{移項}{い|こう}すると（2-14）、
\bkthriline{$(a-a')x+(b-b')=0$\quad がどんな $x$ でも\R{成}{な}り\R{立}{た}つ}

\noindent\begin{minipage}[c]{0.6\linewidth}
\textbf{\R{数}{かず}を\R{入}{い}れて}\par
どんな $x$ でもよいので、\R{計算}{けい|さん}の\R{楽}{らく}な\R{数}{かず}を\R{入}{い}れます。$x=0$ を\R{入}{い}れると、$(a-a')\times0+(b-b')=0$ なので $b-b'=0$、つまり $b=b'$ です。すると\R{式}{しき}は $(a-a')x=0$ になります。ここに $x=1$ を\R{入}{い}れると $a-a'=0$、つまり $a=a'$ です。

\textbf{グラフで}\par
$y=(a-a')x+(b-b')$ は、\R{傾}{かたむ}き $a-a'$、\R{切片}{せっ|ぺん} $b-b'$ の\R{直線}{ちょく|せん}です（3-5）。それがどこでも $y=0$ なら、\R{直線}{ちょく|せん}は $x$ \R{軸}{じく}そのものです。$x$ \R{軸}{じく}は\R{傾}{かたむ}きが0、\R{切片}{せっ|ぺん}も0なので、$a-a'=0$、$b-b'=0$ です。
\end{minipage}\hfill
\begin{minipage}[c]{0.37\linewidth}\centering
\begin{tikzpicture}
\begin{axis}[bkplot, width=46mm, height=46mm, xmin=-2.6, xmax=3, ymin=-2.6, ymax=3, xtick={-2,2}, ytick={-2,2}, xlabel={$x$}, ylabel={$y$}]
  \addplot[bkGray, line width=0.9pt, dashed, domain=-2.5:1.5] {1.2*x+0.8};
  \addplot[bkGray, line width=0.9pt, dashed, domain=-2.4:2.8] {-0.5*x-1};
  \addplot[bkRed, line width=2.4pt, domain=-2.5:2.8] {0};
  \node[font=\sffamily\scriptsize, text=bkRed!80!black, anchor=south west] at (axis cs:1.2,0.1) {$y=0$};
\end{axis}
\end{tikzpicture}
\bkfigcaption{どこでも $y=0$ になる\R{直線}{ちょく|せん}は、$x$ \R{軸}{じく}（\R{赤}{あか}）だけ。\R{点線}{てん|せん}は\R{傾}{かたむ}きや\R{切片}{せっ|ぺん}が0でない\R{直線}{ちょく|せん}}
\end{minipage}

\vspace{3pt}
2\R{次式}{じ|しき}でも\R{同}{おな}じです。$Ax^2+Bx+C=0$ がどんな $x$ でも\R{成}{な}り\R{立}{た}つとき、$x=0$ で $C=0$。$x=1$ で $A+B=0$、$x=-1$ で $A-B=0$。この2つをたすと $2A=0$ なので $A=0$、すると $B=0$ です（3-6の\R{加減法}{か|げん|ほう}）。だから、\R{両辺}{りょう|へん}の\R{差}{さ}の\R{係数}{けい|すう}はどれも0で、\R{係数}{けい|すう}はそれぞれ\R{等}{ひと}しくなります。

\begin{column}{\R{方程式}{ほう|てい|しき}では、\R{係数}{けい|すう}を\R{比}{くら}べてはいけない}
$3x+1=x+5$ は\R{方程式}{ほう|てい|しき}で、$x=2$ のときだけ\R{成}{な}り\R{立}{た}ちます（$3\times2+1=7$、$2+5=7$）。もし\R{係数}{けい|すう}を\R{比}{くら}べると「$3=1$」となって、おかしなことになります。\R{係数}{けい|すう}を\R{比}{くら}べてよいのは、どんな $x$ でも\R{成}{な}り\R{立}{た}つ\R{恒等式}{こう|とう|しき}のときだけです。\R{問題}{もん|だい}に「\R{恒等式}{こう|とう|しき}になるように」「どんな $x$ でも\R{成}{な}り\R{立}{た}つように」と\R{書}{か}いてあるかを、まず\R{確}{たし}かめましょう。
\end{column}

\bksubsection{9}{\R{数}{かず}を\R{代入}{だい|にゅう}して\R{係数}{けい|すう}を\R{決}{き}める}

\leavevmode\label{bk:3-8:dainyu}\R{恒等式}{こう|とう|しき}なら、どんな $x$ を\R{入}{い}れても\R{成}{な}り\R{立}{た}ちます。この\R{性質}{せい|しつ}を、わからない\R{数}{かず}を\R{決}{き}めるのに\R{使}{つか}う\R{方法}{ほう|ほう}もあります。

\bkthriwaku{$5x+1=A(x+1)+B(x-1)$ が\R{恒等式}{こう|とう|しき}になるような、\R{数}{かず} $A$ と $B$ は？}

\bkthrihead{\R{方法}{ほう|ほう}1：\R{都合}{つ|ごう}のよい $x$ を\R{入}{い}れる}
\R{右辺}{う|へん}には $A$ と $B$ の2つがあって、このままでは1つに\R{決}{き}まりません。そこで、どちらか1つが\R{消}{き}えるような $x$ を\R{選}{えら}びます。

\noindent\textbf{$x=1$ を\R{入}{い}れる}\quad $x-1=0$ になるので、$B$ の\R{項}{こう}が\R{消}{き}えます。
\[ 5\times1+1=A(1+1)+B(1-1) \]
\R{左}{ひだり}は6、\R{右}{みぎ}は $2A+B\times0=2A$ なので、$6=2A$。だから $A=3$ です。

\noindent\textbf{$x=-1$ を\R{入}{い}れる}\quad $x+1=0$ になるので、$A$ の\R{項}{こう}が\R{消}{き}えます。
\[ 5\times(-1)+1=A(-1+1)+B(-1-1) \]
\R{左}{ひだり}は $-4$、\R{右}{みぎ}は $A\times0-2B=-2B$ なので、$-4=-2B$。だから $B=2$ です。

$A=3$、$B=2$ です。\R{確}{たし}かめに、\R{選}{えら}ばなかった $x=2$ を\R{入}{い}れると、\R{左}{ひだり}は $5\times2+1=11$、\R{右}{みぎ}は $3\times3+2\times1=11$ で\R{合}{あ}います。

\R{好}{す}きな $x$ を\R{入}{い}れてよいのは、これが\textbf{\R{恒等式}{こう|とう|しき}だから}です。$x=1$ でも $x=-1$ でも $x=2$ でも、\R{同}{おな}じ $A$・$B$ で\R{成}{な}り\R{立}{た}たなければなりません。だから、\R{都合}{つ|ごう}のよい $x$ で\R{求}{もと}めた $A$・$B$ が、そのまま\R{答}{こた}えになります。

\bkthrihead{\R{方法}{ほう|ほう}2：\R{展開}{てん|かい}して\R{係数}{けい|すう}を\R{比}{くら}べる}
\R{右辺}{う|へん}を\R{展開}{てん|かい}して、$x$ の\R{項}{こう}と\R{数}{かず}だけの\R{項}{こう}に\R{整理}{せい|り}します。
\begin{align*}
A(x+1)+B(x-1)&=Ax+A+Bx-B\\
&=(A+B)x+(A-B)
\end{align*}
\R{左辺}{さ|へん} $5x+1$ と\R{係数}{けい|すう}を\R{比}{くら}べると、$A+B=5$、$A-B=1$ です。これは\R{連立方程式}{れん|りつ|ほう|てい|しき}です（3-6）。2つの\R{式}{しき}をたすと $2A=6$ で $A=3$。$A+B=5$ に\R{入}{い}れて $B=2$ です。

\noindent\begin{minipage}[c]{0.52\linewidth}
2つの\R{方法}{ほう|ほう}で、\R{同}{おな}じ $A=3$、$B=2$ が\R{出}{で}ました。\R{方法}{ほう|ほう}1は\R{計算}{けい|さん}が\R{短}{みじか}く、\R{方法}{ほう|ほう}2は「なぜその\R{答}{こた}えか」がよく\R{見}{み}えます。どちらで\R{求}{もと}めても、もう\R{片方}{かた|ほう}で\R{確}{たし}かめると\R{安心}{あん|しん}です。
\end{minipage}\hfill
\begin{minipage}[c]{0.45\linewidth}\centering
\begin{tikzpicture}[x=1mm, y=1mm]
  \node[bk node, minimum width=22mm, font=\sffamily\small] (q) at (24,14) {$5x+1=$\\$A(x+1)+B(x-1)$};
  \node[bk node, minimum width=18mm, font=\sffamily\footnotesize, draw=bkBlue] (m1) at (8,-7) {\R{方法}{ほう|ほう}1\\$x=\pm1$ を\R{入}{い}れる};
  \node[bk node, minimum width=18mm, font=\sffamily\footnotesize, draw=bkGreen] (m2) at (40,-7) {\R{方法}{ほう|ほう}2\\\R{係数}{けい|すう}を\R{比}{くら}べる};
  \node[font=\sffamily\bfseries\small, text=bkRed!80!black] (ans) at (24,-27) {$A=3$、$B=2$};
  \draw[bk link] (q) -- (m1); \draw[bk link] (q) -- (m2);
  \draw[bk link] (m1) -- (ans); \draw[bk link] (m2) -- (ans);
\end{tikzpicture}
\end{minipage}

\vspace{3pt}
この\R{形}{かたち}は、\R{分数}{ぶん|すう}の\R{式}{しき}を\R{分}{わ}ける\R{計算}{けい|さん}（\R{部分分数分解}{ぶ|ぶん|ぶん|すう|ぶん|かい}）で、そのまま\R{使}{つか}います。この\R{章}{しょう}の\R{終}{お}わりのほう（\pageref{bk:3-8:bubun}ページ）で、やってみます。

\begin{seikaku}
\textbf{なぜ、\R{都合}{つ|ごう}のよい2つの $x$ だけで\R{決}{き}めてよいの？}\par
①\enspace $5x+1$ も $A(x+1)+B(x-1)$ も1\R{次式}{じ|しき}です。1\R{次式}{じ|しき}どうしは、ちがう2つの $x$ で\R{値}{あたい}が\R{合}{あ}えば、すべての $x$ で\R{合}{あ}います。2つの\R{差}{さ}のグラフは\R{直線}{ちょく|せん}で、2\R{点}{てん}で0になる\R{直線}{ちょく|せん}は $x$ \R{軸}{じく}しかないからです（\R{道具}{どう|ぐ}8のわけ）。\par
②\enspace \R{部分分数分解}{ぶ|ぶん|ぶん|すう|ぶん|かい}（\pageref{bk:3-8:bubun}ページ）では、はじめに\R{両辺}{りょう|へん}に\R{分母}{ぶん|ぼ}をかけて（\R{分母}{ぶん|ぼ}をはらって）、\R{分数}{ぶん|すう}のない\R{恒等式}{こう|とう|しき}にしてから $x$ を\R{代入}{だい|にゅう}します。だから、もとの\R{分数}{ぶん|すう}の\R{分母}{ぶん|ぼ}を0にする $x$（この\R{例}{れい}では $x=1$ や $x=-1$）を\R{入}{い}れてもかまいません。
\end{seikaku}

\bksubsection{10}{\R{平方完成}{へい|ほう|かん|せい}：\R{正方形}{せい|ほう|けい}をつくる}

\leavevmode\label{bk:3-8:heihou}\R{道具}{どう|ぐ}8で、$x^2+6x+10=(x+3)^2+1$ になることを、\R{係数}{けい|すう}を\R{比}{くら}べて\R{求}{もと}めました。このように、2\R{次式}{じ|しき}を「（かっこ）の2\R{乗}{じょう}＋\R{数}{かず}」の\R{形}{かたち}に\R{直}{なお}すことを\textbf{\R{平方完成}{へい|ほう|かん|せい}}といいます（\R{平方}{へい|ほう}は2\R{乗}{じょう}のこと）。\bkidx{へいほうかんせい}{\R{平方完成}{へい|ほう|かん|せい}}こんどは\R{面積}{めん|せき}の\R{図}{ず}で、なぜ $p=3$ になるのかを\R{見}{み}てみましょう。

\bkthristep{1}{$x^2+6x$ を\R{図}{ず}にする}
$x^2$ は1\R{辺}{ぺん} $x$ の\R{正方形}{せい|ほう|けい}、$6x$ はたて $x$・\R{横}{よこ}6の\R{長方形}{ちょう|ほう|けい}です（\R{図}{ず}\ref{bk:3-8:fig-hei}の\R{左}{ひだり}）。

\bkthristep{2}{\R{長方形}{ちょう|ほう|けい}を\R{半分}{はん|ぶん}に\R{切}{き}って、2\R{辺}{へん}にはりつける}
$6x$ の\R{長方形}{ちょう|ほう|けい}を、\R{横}{よこ}3ずつの2まい（$3x$ と $3x$）に\R{切}{き}り、1まいを\R{正方形}{せい|ほう|けい}の\R{右}{みぎ}に、もう1まいを\R{下}{した}にはりつけます（\R{図}{ず}\ref{bk:3-8:fig-hei}のまん\R{中}{なか}）。

\bkthristep{3}{\R{欠}{か}けた\R{角}{かど}をうめる}
\R{右下}{みぎ|した}に、1\R{辺}{ぺん}3の\R{正方形}{せい|ほう|けい}の\R{形}{かたち}の\R{欠}{か}けがあります。ここに $3\times3=9$ をうめると、1\R{辺}{ぺん} $x+3$ の\R{大}{おお}きな\R{正方形}{せい|ほう|けい}が\R{完成}{かん|せい}します（\R{図}{ず}\ref{bk:3-8:fig-hei}の\R{右}{みぎ}）。うめた9は、\R{本当}{ほん|とう}はないものなので、あとで\R{返}{かえ}します。

\begin{bkfigure}
\begin{tikzpicture}[x=1mm, y=1mm]
  % 左：x^2 と 6x  % noruby
  \fill[bkBlue!14] (0,0) rectangle (20,20);
  \fill[bkOrange!25] (20,0) rectangle (32,20);
  \draw[bkInk, line width=0.8pt] (0,0) rectangle (32,20);
  \draw[bkInk, line width=0.6pt] (20,0) -- (20,20);
  \draw[bkInk, line width=0.4pt, dashed] (26,0) -- (26,20);
  \node[font=\small] at (10,10) {$x^2$};
  \node[font=\small] at (26,10) {$6x$};
  \bkbrace{(0,21.5)}{(20,21.5)}{$x$}{1}\bkbrace{(20,21.5)}{(32,21.5)}{6}{1}
  \draw[bk arrow, bkGray] (36,10) -- (44,10);
  % 中：3x を右と下に  % noruby
  \begin{scope}[shift={(50,-3)}]
  \fill[bkBlue!14] (0,6) rectangle (20,26);
  \fill[bkOrange!25] (20,6) rectangle (26,26);
  \fill[bkOrange!25] (0,0) rectangle (20,6);
  \draw[bkRed!80!black, line width=0.7pt, dashed] (20,0) rectangle (26,6);
  \draw[bkInk, line width=0.8pt] (0,0) -- (20,0) -- (20,6) -- (26,6) -- (26,26) -- (0,26) -- cycle;
  \draw[bkInk, line width=0.5pt] (20,6) -- (20,26) (0,6) -- (20,6);
  \node[font=\small] at (10,16) {$x^2$};
  \node[font=\footnotesize] at (23,16) {\rotatebox{90}{$3x$}};
  \node[font=\footnotesize] at (10,3) {$3x$};
  \node[font=\footnotesize, text=bkRed!80!black] at (23,3) {？};
  \bkbrace{(0,27.5)}{(20,27.5)}{$x$}{1}\bkbrace{(20,27.5)}{(26,27.5)}{3}{1}
  \end{scope}
  \draw[bk arrow, bkGray] (80,10) -- (88,10);
  % 右：9 をうめて (x+3)^2  % noruby
  \begin{scope}[shift={(94,-3)}]
  \fill[bkBlue!14] (0,6) rectangle (20,26);
  \fill[bkOrange!25] (20,6) rectangle (26,26);
  \fill[bkOrange!25] (0,0) rectangle (20,6);
  \fill[bkRed!35] (20,0) rectangle (26,6);
  \draw[bkInk, line width=0.8pt] (0,0) rectangle (26,26);
  \draw[bkInk, line width=0.5pt] (20,0) -- (20,26) (0,6) -- (26,6);
  \node[font=\small] at (10,16) {$x^2$};
  \node[font=\footnotesize] at (23,3) {9};
  \bkbrace{(0,27.5)}{(26,27.5)}{$x+3$}{1}
  \node[font=\small, anchor=west] at (27.5,13) {$x+3$};
  \end{scope}
\end{tikzpicture}
\bkfigcaption{$x^2+6x$ に9をたすと、1\R{辺}{ぺん} $x+3$ の\R{正方形}{せい|ほう|けい}になる}\label{bk:3-8:fig-hei}
\end{bkfigure}

\R{図}{ず}から、$x^2+6x+9=(x+3)^2$ です。だから $x^2+6x=(x+3)^2-9$。これに10をたすと、
\begin{align*}
x^2+6x+10&=(x+3)^2-9+10\\
&=(x+3)^2+1
\end{align*}
\R{係数}{けい|すう}を\R{比}{くら}べた\R{答}{こた}え（\R{道具}{どう|ぐ}8）と\R{同}{おな}じになりました。3は「$x$ の\R{係数}{けい|すう}6の\R{半分}{はん|ぶん}」で、9はその2\R{乗}{じょう}です。\R{長方形}{ちょう|ほう|けい}を\R{半分}{はん|ぶん}に\R{切}{き}って2\R{辺}{へん}にはったから\R{半分}{はん|ぶん}、\R{欠}{か}けた\R{角}{かど}が\R{正方形}{せい|ほう|けい}だから2\R{乗}{じょう}なのです。

\begin{taisetsu}
\R{平方完成}{へい|ほう|かん|せい}の\R{手順}{て|じゅん}：$x^2+bx+c$ は、$x$ の\R{係数}{けい|すう} $b$ の\R{半分}{はん|ぶん}を $p$ として、\par\vspace{2pt}
$x^2+bx+c=(x+p)^2-p^2+c$\par\vspace{2pt}
（$(x+p)^2$ を\R{作}{つく}るためにたした $p^2$ を、すぐにひいて\R{返}{かえ}す）
\end{taisetsu}

\bkthrihead{\R{平方完成}{へい|ほう|かん|せい}すると、いちばん\R{小}{ちい}さい\R{値}{あたい}が\R{見}{み}える}
\noindent\begin{minipage}[c]{0.53\linewidth}
$(x+3)^2$ は2\R{乗}{じょう}なので、0かプラスです。マイナスの\R{数}{かず}も2\R{乗}{じょう}すればプラスになるからです（2-2）。だから $(x+3)^2+1$ は、いつも1\R{以上}{い|じょう}です。ちょうど1になるのは、$(x+3)^2=0$、つまり $x=-3$ のときだけです。

\R{表}{ひょう}とグラフ（\R{赤}{あか}い\R{点}{てん}が $(-3,\ 1)$）でも、$x=-3$ でいちばん\R{小}{ちい}さな1になり、そこから\R{左右}{さ|ゆう}に\R{同}{おな}じように\R{大}{おお}きくなっています。
\end{minipage}\hfill
\begin{minipage}[c]{0.44\linewidth}\centering
\begin{tikzpicture}
\begin{axis}[bkplot, width=52mm, height=50mm, xmin=-6.6, xmax=0.9, ymin=0, ymax=11.5, xtick={-6,-3}, ytick={1,5,10}, xlabel={$x$}, ylabel={$y$}]
  \addplot[bkcurve, domain=-6.2:0.1, samples=100] {x^2+6*x+10};
  \addplot[only marks, mark=*, mark size=1.6pt, bkInk] coordinates {(-6,10) (-5,5) (-4,2) (-2,2) (-1,5) (0,10)};
  \addplot[only marks, mark=*, mark size=2.2pt, bkRed] coordinates {(-3,1)};
\end{axis}
\end{tikzpicture}
{\small
\begin{bktable}{colspec={Q[c,m,wd=8mm] *{7}{Q[c,m,wd=4mm]}}, rowsep=0.5pt, colsep=1.6pt, cells={font=\scriptsize}}
$x$ & $-6$ & $-5$ & $-4$ & $-3$ & $-2$ & $-1$ & 0 \\
$y$ & 10 & 5 & 2 & 1 & 2 & 5 & 10 \\
\end{bktable}}
\end{minipage}

\vspace{3pt}
\R{表}{ひょう}を\R{作}{つく}って\R{調}{しら}べるだけでは、$-3.1$ や $-2.9$ のほうが\R{小}{ちい}さいかもしれない、という\R{心配}{しん|ぱい}がのこります。\R{平方完成}{へい|ほう|かん|せい}した\R{式}{しき}なら、どの $x$ でも1より\R{小}{ちい}さくならないことが、はっきり\R{言}{い}えます。

\bkthrihead{$x^2$ の\R{係数}{けい|すう}が1でないとき}
$2x^2-8x+5$ のように $x^2$ に\R{数}{かず}がかかっているときは、はじめに $x^2$ と $x$ の\R{項}{こう}を、その\R{数}{かず}でくくります（\R{道具}{どう|ぐ}6）。
\begin{enumerate}[label={\bkcircled[bkOrange]{\arabic*}}, leftmargin=2.4\zw, itemsep=2pt, topsep=3pt]
\item 2でくくる：\quad $2x^2-8x+5=2(x^2-4x)+5$
\item かっこの\R{中}{なか}を\R{平方完成}{へい|ほう|かん|せい}する（$-4$ の\R{半分}{はん|ぶん}は $-2$）：\quad $x^2-4x=(x-2)^2-4$\\
だから\quad $2(x^2-4x)+5=2\{(x-2)^2-4\}+5$
\item かっこをはずす（\R{分配法則}{ぶん|ぱい|ほう|そく}）：\quad $2(x-2)^2-8+5=2(x-2)^2-3$
\end{enumerate}
\R{波}{なみ}かっこ $\{\ \}$ は、かっこの\R{外}{そと}がわにつけるかっこです。$2\times(-4)=-8$ をわすれないようにします。

\bkthrihead{\bkhatenaref の\R{答}{こた}え：ボールのいちばん\R{高}{たか}い\R{点}{てん}}
$h=20t-5t^2$ を、$t^2$ の\R{項}{こう}を\R{先}{さき}にして $-5t^2+20t$ と\R{書}{か}き、$-5$ でくくります（$-5\times t^2=-5t^2$、$-5\times(-4t)=20t$）。
\begin{enumerate}[label={\bkcircled[bkOrange]{\arabic*}}, leftmargin=2.4\zw, itemsep=2pt, topsep=3pt]
\item $-5$ でくくる：\quad $h=-5(t^2-4t)$
\item かっこの\R{中}{なか}を\R{平方完成}{へい|ほう|かん|せい}する：\quad $h=-5\{(t-2)^2-4\}$
\item かっこをはずす（$-5\times(-4)=20$）：\quad $h=-5(t-2)^2+20$
\end{enumerate}
\R{順番}{じゅん|ばん}を\R{入}{い}れかえると、はてな？の $h=20-5(t-2)^2$ です。\R{展開}{てん|かい}して\R{確}{たし}かめると、
\[ 20-5(t^2-4t+4)=20t-5t^2 \]
で、たしかに\R{同}{おな}じ\R{式}{しき}（\R{恒等式}{こう|とう|しき}）です。

\noindent\begin{minipage}[c]{0.55\linewidth}
$(t-2)^2$ は0かプラスなので、$5(t-2)^2$ も0かプラスです。20からそれをひいた $h$ は、20より\R{大}{おお}きくなれません。$h$ がちょうど20になるのは $t=2$ のときだけです。だから、\textbf{いちばん\R{高}{たか}いのは2\R{秒後}{びょう|ご}で、\R{高}{たか}さ\qty{20}{m}}です。1.9\R{秒}{びょう}も2.1\R{秒}{びょう}も、$20-5\times0.1^2=19.95$ で、20より\R{低}{ひく}くなります。
\end{minipage}\hfill
\begin{minipage}[c]{0.42\linewidth}\centering
\begin{tikzpicture}
\begin{axis}[bkplot, width=52mm, height=48mm, xmin=0, xmax=4.6, ymin=0, ymax=25, xtick={1,2,3,4}, ytick={5,10,15,20}, xlabel={$t$（\R{秒}{びょう}）}, ylabel={$h$（m）}]
  \addplot[bkcurve, domain=0:4] {20*x-5*x^2};
  \draw[bkRed, line width=0.7pt, dash pattern=on 2pt off 1.5pt] (axis cs:2,0) -- (axis cs:2,20) -- (axis cs:0,20);
  \addplot[only marks, mark=*, mark size=2.2pt, bkRed] coordinates {(2,20)};
  \node[font=\sffamily\scriptsize, text=bkRed!80!black, anchor=south west] at (axis cs:2.1,20.3) {いちばん\R{高}{たか}い};
\end{axis}
\end{tikzpicture}
\end{minipage}

\vspace{3pt}
このように、\R{式}{しき}の\R{形}{かたち}を\R{変}{か}えると、もとの\R{形}{かたち}では\R{見}{み}えなかったことが\R{見}{み}えてきます。グラフのいちばん\R{高}{たか}い\R{点}{てん}・\R{低}{ひく}い\R{点}{てん}のことは、3-9でくわしく\R{学}{まな}びます。5-7では、\R{同}{おな}じ\R{答}{こた}えを\R{微分}{び|ぶん}で\R{求}{もと}めます。

\Needspace{12\baselineskip}
\begin{motto}[3]
\textbf{3-16への\R{予告}{よ|こく}：ばねの\R{実験}{じっ|けん}に、いちばんよい\R{直線}{ちょく|せん}を\R{引}{ひ}く}\par
おもりを $x$ \R{個}{こ}つるしたときのばねののび $y$〔cm〕をはかると、$(1,\ 2.0)$、$(2,\ 4.3)$、$(3,\ 5.8)$ でした。\R{点}{てん}は\R{直線}{ちょく|せん}にぴったりはのりません（2-6のずれ）。そこで、\R{直線}{ちょく|せん} $y=kx$ からのずれ $y-kx$ を2\R{乗}{じょう}してたした $S$ が、いちばん\R{小}{ちい}さくなる $k$ を\R{選}{えら}びます（なぜ2\R{乗}{じょう}するかは3-16で）。
\[ S=(2.0-k)^2+(4.3-2k)^2+(5.8-3k)^2 \]
1つずつ $(a-b)^2=a^2-2ab+b^2$ で\R{展開}{てん|かい}します。
\begin{align*}
(2.0-k)^2&=4-4k+k^2\\
(4.3-2k)^2&=18.49-17.2k+4k^2\\
(5.8-3k)^2&=33.64-34.8k+9k^2
\end{align*}
\R{同類項}{どう|るい|こう}をまとめると $S=14k^2-56k+56.13$。14でくくって\R{平方完成}{へい|ほう|かん|せい}すると、
\[ S=14(k^2-4k)+56.13=14(k-2)^2+0.13 \]
$14(k-2)^2$ は0かプラスなので、$S$ がいちばん\R{小}{ちい}さくなるのは $k=2$ のときです。だから、このばねは「おもり1\R{個}{こ}で\qty{2}{cm}のびる」と\R{決}{き}められます。\R{実験}{じっ|けん}からばね\R{定数}{てい|すう}を\R{決}{き}めるこの\R{方法}{ほう|ほう}（\R{最小二乗法}{さい|しょう|に|じょう|ほう}）は、3-16でくわしく\R{学}{まな}びます。
\end{motto}

\Needspace{12\baselineskip}
\begin{motto}[3]
\textbf{\R{係数}{けい|すう}のならび：1・2・1、1・3・3・1}\par
\noindent\begin{minipage}[c]{0.6\linewidth}
$(a+b)^2=a^2+2ab+b^2$ の\R{係数}{けい|すう}は1・2・1、$(a+b)^3=a^3+3a^2b+3ab^2+b^3$ は1・3・3・1です。\R{右}{みぎ}のように、\R{上}{うえ}の2つの\R{数}{かず}をたして\R{下}{した}の\R{数}{かず}を\R{作}{つく}ると、$(a+b)^4$ の\R{係数}{けい|すう}1・4・6・4・1も\R{作}{つく}れます。この\R{三角形}{さん|かく|けい}を\textbf{パスカルの\R{三角形}{さん|かく|けい}}といい、3-17で、コインの\R{表}{おもて}・\R{裏}{うら}の\R{出方}{で|かた}を\R{数}{かぞ}えるのに\R{使}{つか}います。
\end{minipage}\hfill
\begin{minipage}[c]{0.36\linewidth}\centering
\begin{tikzpicture}[x=1mm, y=1mm]
  \foreach \r/\row in {0/{1}, 1/{1,1}, 2/{1,2,1}, 3/{1,3,3,1}, 4/{1,4,6,4,1}} {
    \foreach \v [count=\i from 0] in \row {\node[font=\small, inner sep=1pt] at ({(\i-\r/2)*8},{-\r*9}) {\v};}}
  \draw[bkOrange!80!black, line width=0.6pt] (-7,-20.5) -- (-5,-24.5) (-1,-20.5) -- (-3,-24.5);
\end{tikzpicture}
\end{minipage}
\par\vspace{4pt}
\textbf{3つの\R{和}{わ}の2\R{乗}{じょう}}\par
$(a+b+c)^2$ は、1\R{辺}{ぺん} $a+b+c$ の\R{正方形}{せい|ほう|けい}を9つに\R{分}{わ}けて\R{考}{かんが}えます。$a^2$、$b^2$、$c^2$ が1まいずつ、$ab$、$bc$、$ca$ が2まいずつで、
\[ (a+b+c)^2=a^2+b^2+c^2+2ab+2bc+2ca \]
\end{motto}

\begin{kotoba}
\small
\bkword{\R{単項式}{たん|こう|しき}・\R{多項式}{た|こう|しき}}{\R{数}{かず}と\R{文字}{も|じ}のかけ\R{算}{ざん}だけの\R{式}{しき}が\R{単項式}{たん|こう|しき}。いくつかの\R{項}{こう}の\R{和}{わ}が\R{多項式}{た|こう|しき}。}
\bkword{\R{次数}{じ|すう}}{かけ\R{合}{あ}わせた\R{文字}{も|じ}の\R{個数}{こ|すう}。$3a^2b$ は3\R{次}{じ}。\R{多項式}{た|こう|しき}はいちばん\R{大}{おお}きい\R{項}{こう}の\R{次数}{じ|すう}。}
\bkword{2\R{次式}{じ|しき}}{いちばん\R{大}{おお}きい\R{次数}{じ|すう}が2の\R{多項式}{た|こう|しき}。$x^2+6x+10$ など。}
\bkword{\R{展開}{てん|かい}}{かっこのある\R{式}{しき}をかけ\R{算}{ざん}して、かっこのない\R{多項式}{た|こう|しき}にすること。}
\bkword{\R{共通因数}{きょう|つう|いん|すう}・くくる}{どの\R{項}{こう}にもかかっている\R{数}{かず}や\R{文字}{も|じ}。それを\R{外}{そと}に\R{出}{だ}すことが「くくる」。}
\bkword{\R{恒等式}{こう|とう|しき}}{\R{文字}{も|じ}にどんな\R{数}{かず}を\R{入}{い}れても\R{成}{な}り\R{立}{た}つ\R{等式}{とう|しき}。}
\bkword{\R{係数}{けい|すう}を\R{比}{くら}べる}{\R{恒等式}{こう|とう|しき}の\R{両辺}{りょう|へん}で、\R{同}{おな}じ\R{次数}{じ|すう}の\R{係数}{けい|すう}を\R{等}{ひと}しいとおいて、わからない\R{数}{かず}を\R{求}{もと}めること。}
\bkword{\R{平方完成}{へい|ほう|かん|せい}}{2\R{次式}{じ|しき}を「（かっこ）の2\R{乗}{じょう}＋\R{数}{かず}」の\R{形}{かたち}にすること。}
\end{kotoba}

%% ---------------- 例題 ----------------
\bksection{reidai}{\R{例題}{れい|だい}でたしかめよう}

\begin{reidai}[\R{単項式}{たん|こう|しき}のかけ\R{算}{ざん}]
\R{次}{つぎ}の\R{計算}{けい|さん}をしましょう。\R{答}{こた}えの\R{次数}{じ|すう}も\R{答}{こた}えましょう。\par
（1）$3a^2b\times2ab$\qquad（2）$(-4x)\times5x^2$\qquad（3）$(2x)^2$\par
\kangaekata ×を\R{全部}{ぜん|ぶ}\R{書}{か}いて、\R{数}{かず}は\R{数}{かず}どうし、\R{同}{おな}じ\R{文字}{も|じ}は\R{同}{おな}じ\R{文字}{も|じ}どうしでまとめます。\R{符号}{ふ|ごう}は\R{数}{かず}といっしょに\R{計算}{けい|さん}します。\par
\shiki（1）\nobreak $3\times2\times a\times a\times a\times b\times b=6a^3b^2$\par
\hspace*{2.6\zw}（2）\nobreak $(-4)\times5\times x\times x\times x=-20x^3$\par
\hspace*{2.6\zw}（3）\nobreak $2x\times2x=2\times2\times x\times x=4x^2$\par
\kotae（1）\nobreak $6a^3b^2$、5\R{次}{じ}\qquad（2）\nobreak $-20x^3$、3\R{次}{じ}\qquad（3）\nobreak $4x^2$、2\R{次}{じ}
\end{reidai}

\begin{reidai}[かっこ×かっこの\R{展開}{てん|かい}]
\R{次}{つぎ}の\R{式}{しき}を\R{展開}{てん|かい}しましょう。\par
（1）$(x+3)(y+2)$\qquad（2）$(2x+1)(x-3)$\par
\kangaekata \R{全部}{ぜん|ぶ}の\R{組}{く}み\R{合}{あ}わせ（$2\times2=4$ \R{組}{くみ}）をかけてたします。（2）は $x-3=x+(-3)$ と\R{見}{み}て、\R{最後}{さい|ご}に\R{同類項}{どう|るい|こう}をまとめます。\par
\shiki（1）\nobreak $x\times y+x\times2+3\times y+3\times2=xy+2x+3y+6$\par
\hspace*{2.6\zw}（2）\nobreak $2x\times x+2x\times(-3)+1\times x+1\times(-3)$\par
\hspace*{2.6\zw}\phantom{（2）}$=2x^2-6x+x-3=2x^2-5x-3$\par
\kotae（1）\nobreak $xy+2x+3y+6$\qquad（2）\nobreak $2x^2-5x-3$
\end{reidai}

\begingroup\tcbset{equal height group=bkthrijA}%
\noindent\begin{minipage}[t]{0.49\linewidth}
\begin{jibun}
\R{計算}{けい|さん}をし、\R{次数}{じ|すう}も\R{答}{こた}えましょう。\par
（1）$4xy\times3x$\par
（2）$(-2a)\times3ab$\par
（3）$(3x)^2$
\bkjans{（1）\nobreak $12x^2y$、3\R{次}{じ}。（2）\nobreak $-6a^2b$、3\R{次}{じ}。（3）\nobreak $9x^2$、2\R{次}{じ}}
\end{jibun}
\end{minipage}\hfill
\begin{minipage}[t]{0.49\linewidth}
\begin{jibun}
\R{展開}{てん|かい}しましょう。\par
（1）$(a+4)(b+1)$\par
（2）$(3x-2)(x+4)$
\bkjans{（1）\nobreak $ab+a+4b+4$。（2）\nobreak $3x^2+12x-2x-8=3x^2+10x-8$}
\end{jibun}
\end{minipage}\par
\endgroup

\begin{reidai}[よく\R{使}{つか}う\R{展開}{てん|かい}の\R{形}{かたち}]
\R{次}{つぎ}の\R{式}{しき}を\R{展開}{てん|かい}しましょう。\par
（1）$(x+4)^2$\qquad（2）$(12-2x)^2$\qquad（3）$(x+5)(x-5)$\qquad（4）$(x-3)(x+5)$\par
\kangaekata どれも「\R{全部}{ぜん|ぶ}の\R{組}{く}み\R{合}{あ}わせ」でできますが、\R{道具}{どう|ぐ}4の\R{表}{ひょう}の\R{形}{かたち}を\R{使}{つか}うと\R{速}{はや}く\R{書}{か}けます。（2）は $a=12$、$b=2x$ の $(a-b)^2$、（4）は $a=-3$、$b=5$ の $(x+a)(x+b)$ です。\par
\shiki（1）\nobreak $x^2+2\times x\times4+4^2=x^2+8x+16$\par
\hspace*{2.6\zw}（2）\nobreak $12^2-2\times12\times2x+(2x)^2=144-48x+4x^2$\par
\hspace*{2.6\zw}（3）\nobreak $x^2-5^2=x^2-25$\par
\hspace*{2.6\zw}（4）\nobreak \R{和}{わ}は $-3+5=2$、\R{積}{せき}は $(-3)\times5=-15$。だから $x^2+2x-15$\par
\kotae（1）\nobreak $x^2+8x+16$\qquad（2）\nobreak $144-48x+4x^2$\par
\hspace*{2.6\zw}（3）\nobreak $x^2-25$\qquad（4）\nobreak $x^2+2x-15$
\end{reidai}

\begin{reidai}[\R{展開}{てん|かい}で、くふうして\R{計算}{けい|さん}する]
くふうして\R{計算}{けい|さん}しましょう。\par
（1）$102^2$\qquad（2）$98^2$\qquad（3）$101\times99$\qquad（4）$1.1^2$\par
\kangaekata 100や1に\R{近}{ちか}い\R{数}{かず}を、「100＋\R{小}{ちい}さい\R{数}{かず}」「1＋\R{小}{ちい}さい\R{数}{かず}」と\R{見}{み}て、\R{展開}{てん|かい}の\R{形}{かたち}を\R{使}{つか}います。\par
\shiki（1）\nobreak $(100+2)^2=10000+400+4=10404$\par
\hspace*{2.6\zw}（2）\nobreak $(100-2)^2=10000-400+4=9604$\par
\hspace*{2.6\zw}（3）\nobreak $(100+1)(100-1)=10000-1=9999$\par
\hspace*{2.6\zw}（4）\nobreak $(1+0.1)^2=1+0.2+0.01=1.21$\par
\kotae（1）\nobreak 10404\qquad（2）\nobreak 9604\qquad（3）\nobreak 9999\qquad（4）\nobreak 1.21
\end{reidai}

\begingroup\tcbset{equal height group=bkthrijB}%
\noindent\begin{minipage}[t]{0.49\linewidth}
\begin{jibun}
\R{展開}{てん|かい}しましょう。\par
（1）$(x+6)^2$\quad（2）$(10-2x)^2$\par
（3）$(x+7)(x-7)$\par
（4）$(x+4)(x-6)$
\bkjans{（1）\nobreak $x^2+12x+36$。（2）\nobreak $100-40x+4x^2$。（3）\nobreak $x^2-49$。（4）\nobreak $x^2-2x-24$}
\end{jibun}
\end{minipage}\hfill
\begin{minipage}[t]{0.49\linewidth}
\begin{jibun}
くふうして\R{計算}{けい|さん}しましょう。\par
（1）$103^2$\quad（2）$99^2$\par
（3）$52\times48$\quad（4）$1.2^2$
\bkjans{（1）\nobreak $(100+3)^2=10609$。（2）\nobreak $(100-1)^2=9801$。（3）\nobreak $(50+2)(50-2)=2500-4=2496$。（4）\nobreak $(1+0.2)^2=1+0.4+0.04=1.44$}
\end{jibun}
\end{minipage}\par
\endgroup

\begin{reidai}[\R{共通因数}{きょう|つう|いん|すう}でくくる]
\R{次}{つぎ}の\R{式}{しき}を、\R{共通因数}{きょう|つう|いん|すう}でくくりましょう。\par
（1）$12x^2-96x+144$\qquad（2）$2x^2-8x$\qquad（3）$3x^2-3$\par
\kangaekata どの\R{項}{こう}にもかかっている\R{数}{かず}や\R{文字}{も|じ}を\R{見}{み}つけ、\R{項}{こう}ごとに「\R{共通因数}{きょう|つう|いん|すう}×\R{何}{なに}か」の\R{形}{かたち}にしてから\R{外}{そと}に\R{出}{だ}します。\R{最後}{さい|ご}に\R{展開}{てん|かい}して\R{確}{たし}かめます。\par
\shiki（1）\nobreak $12\times x^2-12\times8x+12\times12=12(x^2-8x+12)$\par
\hspace*{2.6\zw}（2）\nobreak $2x\times x-2x\times4=2x(x-4)$\par
\hspace*{2.6\zw}（3）\nobreak $3\times x^2-3\times1=3(x^2-1)$\par
\kotae（1）\nobreak $12(x^2-8x+12)$\qquad（2）\nobreak $2x(x-4)$\qquad（3）\nobreak $3(x^2-1)$
\end{reidai}

\begin{reidai}[\R{恒等式}{こう|とう|しき}か、\R{方程式}{ほう|てい|しき}か]
\R{次}{つぎ}の\R{等式}{とう|しき}は、\R{恒等式}{こう|とう|しき}ですか、\R{方程式}{ほう|てい|しき}ですか。\R{方程式}{ほう|てい|しき}なら、\R{成}{な}り\R{立}{た}つ $x$ も\R{求}{もと}めましょう。\par
（1）$(x+2)^2=x^2+4x+4$\qquad（2）$(x+2)^2=x^2+4$\par
\kangaekata はじめに $x=0$、1 などを\R{入}{い}れて、\R{合}{あ}わない $x$ がないか\R{調}{しら}べます。そのあと\R{左辺}{さ|へん}を\R{展開}{てん|かい}して、\R{右辺}{う|へん}とくらべます。\par
\shiki（1）\nobreak \R{左辺}{さ|へん}を\R{展開}{てん|かい}すると $x^2+4x+4$。\R{右辺}{う|へん}とまったく\R{同}{おな}じ\par
\hspace*{2.6\zw}（2）\nobreak $x=1$ なら、\R{左辺}{さ|へん}は9、\R{右辺}{う|へん}は5で、ちがう\par
\hspace*{2.6\zw}\phantom{（2）}\R{左辺}{さ|へん}を\R{展開}{てん|かい}して $x^2+4x+4=x^2+4$\par
\hspace*{2.6\zw}\phantom{（2）}\R{両辺}{りょう|へん}から $x^2+4$ をひいて $4x=0$。だから $x=0$\par
\kotae（1）\nobreak \R{恒等式}{こう|とう|しき}\qquad（2）\nobreak \R{方程式}{ほう|てい|しき}。\R{成}{な}り\R{立}{た}つのは $x=0$ のときだけ
\end{reidai}

\begingroup\tcbset{equal height group=bkthrijC}%
\noindent\begin{minipage}[t]{0.49\linewidth}
\begin{jibun}
\R{共通因数}{きょう|つう|いん|すう}でくくりましょう。\par
（1）$6x^2+18x-24$\par
（2）$5t^2-20t$\quad（3）$4x^2-4$
\bkjans{（1）\nobreak $6(x^2+3x-4)$。（2）\nobreak $5t(t-4)$。（3）\nobreak $4(x^2-1)$}
\end{jibun}
\end{minipage}\hfill
\begin{minipage}[t]{0.49\linewidth}
\begin{jibun}
\R{恒等式}{こう|とう|しき}か、\R{方程式}{ほう|てい|しき}か。\R{方程式}{ほう|てい|しき}なら $x$ も\R{求}{もと}めましょう。\par
（1）$(x-3)^2=x^2-6x+9$\par
（2）$(x-3)^2=x^2-9$
\bkjans{（1）\nobreak \R{恒等式}{こう|とう|しき}。（2）\nobreak \R{方程式}{ほう|てい|しき}。$-6x=-18$ で $x=3$ のときだけ}
\end{jibun}
\end{minipage}\par
\endgroup

\begin{reidai}[\R{係数}{けい|すう}を\R{比}{くら}べて、\R{数}{かず}を\R{決}{き}める]
\R{次}{つぎ}の\R{等式}{とう|しき}が\R{恒等式}{こう|とう|しき}になるように、\R{数}{かず}を\R{決}{き}めましょう。\par
（1）$x^2+8x+20=(x+p)^2+q$\qquad（2）$3x+7=a(x+2)+b$\par
\kangaekata \R{右辺}{う|へん}を\R{展開}{てん|かい}して、\R{同}{おな}じ\R{次数}{じ|すう}の\R{係数}{けい|すう}どうしを\R{等}{ひと}しいとおきます。\par
\shiki（1）\nobreak \R{右辺}{う|へん}を\R{展開}{てん|かい}すると $x^2+2px+p^2+q$\par
\hspace*{2.6\zw}\phantom{（1）}$x$ の\R{係数}{けい|すう}を\R{比}{くら}べて $2p=8$。だから $p=4$\par
\hspace*{2.6\zw}\phantom{（1）}\R{数}{かず}の\R{項}{こう}を\R{比}{くら}べて $p^2+q=20$。$16+q=20$ だから $q=4$\par
\hspace*{2.6\zw}（2）\nobreak \R{右辺}{う|へん}を\R{展開}{てん|かい}すると $ax+2a+b$\par
\hspace*{2.6\zw}\phantom{（2）}$x$ の\R{係数}{けい|すう}を\R{比}{くら}べて $a=3$\par
\hspace*{2.6\zw}\phantom{（2）}\R{数}{かず}の\R{項}{こう}を\R{比}{くら}べて $2a+b=7$。$6+b=7$ だから $b=1$\par
\kotae（1）\nobreak $p=4$、$q=4$\qquad（2）\nobreak $a=3$、$b=1$
\end{reidai}

\begin{reidai}[\R{数}{かず}を\R{代入}{だい|にゅう}して、\R{数}{かず}を\R{決}{き}める]
$x+7=A(x+3)+B(x-1)$ が\R{恒等式}{こう|とう|しき}になるように、\R{数}{かず} $A$、$B$ を\R{決}{き}めましょう。\R{係数}{けい|すう}を\R{比}{くら}べる\R{方法}{ほう|ほう}でも\R{確}{たし}かめましょう。\par
\kangaekata $x=1$ なら $x-1=0$ で $B$ が\R{消}{き}えます。$x=-3$ なら $x+3=0$ で $A$ が\R{消}{き}えます。\par
\shiki $x=1$ を\R{入}{い}れる：$1+7=A\times4+B\times0$\par
\hspace*{2.6\zw}$8=4A$ なので $A=2$\par
\hspace*{2.6\zw}$x=-3$ を\R{入}{い}れる：$-3+7=A\times0+B\times(-4)$\par
\hspace*{2.6\zw}$4=-4B$ なので $B=-1$\par
\hspace*{2.6\zw}\R{確}{たし}かめ：\R{右辺}{う|へん}を\R{展開}{てん|かい}すると $(A+B)x+(3A-B)$\par
\hspace*{2.6\zw}$A+B=2-1=1$、$3A-B=6+1=7$ で、\R{左辺}{さ|へん} $x+7$ と\R{合}{あ}う\par
\kotae $A=2$、$B=-1$
\end{reidai}

\begingroup\tcbset{equal height group=bkthrijD}%
\noindent\begin{minipage}[t]{0.49\linewidth}
\begin{jibun}
\R{恒等式}{こう|とう|しき}になるように\R{決}{き}めましょう。\par
（1）$x^2+10x+30=$\par
\hspace*{3\zw}$(x+p)^2+q$\par
（2）$4x+5=a(x+1)+b$
\bkjans{（1）\nobreak $2p=10$ で $p=5$、$25+q=30$ で $q=5$。（2）\nobreak $a=4$、$a+b=5$ で $b=1$}
\end{jibun}
\end{minipage}\hfill
\begin{minipage}[t]{0.49\linewidth}
\begin{jibun}
\R{次}{つぎ}の\R{式}{しき}が\R{恒等式}{こう|とう|しき}になるように $A$、$B$ を\R{決}{き}め、\R{係数}{けい|すう}を\R{比}{くら}べる\R{方法}{ほう|ほう}でも\R{確}{たし}かめましょう。\par
$6x+4=A(x+2)+B(x-2)$
\bkjans{$x=2$ で $16=4A$、$A=4$。$x=-2$ で $-8=-4B$、$B=2$（$A+B=6$、$2A-2B=4$ で\R{合}{あ}う）}
\end{jibun}
\end{minipage}\par
\endgroup

\begin{reidai}[\R{平方完成}{へい|ほう|かん|せい}して、いちばん\R{小}{ちい}さい\R{値}{あたい}を\R{求}{もと}める]
\R{次}{つぎ}の\R{式}{しき}を\R{平方完成}{へい|ほう|かん|せい}しましょう。（1）は、いちばん\R{小}{ちい}さい\R{値}{あたい}と、そのときの $x$ も\R{求}{もと}めましょう。\par
（1）$x^2+4x+7$\qquad（2）$x^2-6x+5$\qquad（3）$x^2+3x+1$\par
\kangaekata $x$ の\R{係数}{けい|すう}の\R{半分}{はん|ぶん}を $p$ として $(x+p)^2$ を\R{作}{つく}り、たした $p^2$ をひいて\R{返}{かえ}します。（3）は、\R{半分}{はん|ぶん}が\R{分数}{ぶん|すう}になります。\par
\shiki（1）\nobreak 4の\R{半分}{はん|ぶん}は2。$(x+2)^2-4+7=(x+2)^2+3$\par
\hspace*{2.6\zw}\phantom{（1）}$(x+2)^2$ は0かプラスなので、$x=-2$ のときいちばん\R{小}{ちい}さい\par
\hspace*{2.6\zw}（2）\nobreak $-6$ の\R{半分}{はん|ぶん}は $-3$。$(x-3)^2-9+5=(x-3)^2-4$\par
\vspace{5pt}% 第14回 A2：(3) の分数の行と、上下のルビの行がぶつからないよう、前後に余白をとる
\hspace*{2.6\zw}（3）\nobreak 3の\R{半分}{はん|ぶん}は $\dfrac{3}{2}$、その2\R{乗}{じょう}は $\dfrac{9}{4}$\par\vspace{6pt}
\hspace*{2.6\zw}\phantom{（3）}$\Bigl(x+\dfrac{3}{2}\Bigr)^2-\dfrac{9}{4}+1=\Bigl(x+\dfrac{3}{2}\Bigr)^2-\dfrac{5}{4}$\par\vspace{5pt}
\kotae（1）\nobreak $(x+2)^2+3$。$x=-2$ のときいちばん\R{小}{ちい}さく、3\par\vspace{3pt}
\hspace*{2.6\zw}（2）\nobreak $(x-3)^2-4$\qquad（3）\nobreak $\Bigl(x+\dfrac{3}{2}\Bigr)^2-\dfrac{5}{4}$
\end{reidai}

\begin{jibun}
\R{平方完成}{へい|ほう|かん|せい}しましょう。（1）は、いちばん\R{小}{ちい}さい\R{値}{あたい}と、そのときの $x$ も\R{求}{もと}めましょう。\par
（1）$x^2+8x+10$\par
（2）$x^2-10x+30$\par
（3）$x^2+5x+2$
\bkjans{（1）\nobreak $(x+4)^2-6$、$x=-4$ のときいちばん\R{小}{ちい}さく $-6$。（2）\nobreak $(x-5)^2+5$。（3）\nobreak $\bigl(x+\frac{5}{2}\bigr)^2-\frac{17}{4}$}
\end{jibun}

\begin{reidai}[$x^2$ の\R{係数}{けい|すう}が1でない\R{平方完成}{へい|ほう|かん|せい}と、ボールの\R{最高点}{さい|こう|てん}]
（1）$2x^2-8x+5$ を\R{平方完成}{へい|ほう|かん|せい}しましょう。\par
（2）\R{高}{たか}さ\qty{1}{m}の\R{所}{ところ}から\qty{10}{m/s}で\R{真上}{ま|うえ}に\R{投}{な}げたボールの\R{高}{たか}さは、$h=10t-5t^2+1$（m）です。いちばん\R{高}{たか}くなるのは\R{何秒後}{なん|びょう|ご}で、\R{何}{なん}mですか。\par
\kangaekata $x^2$（$t^2$）の\R{係数}{けい|すう}で $x^2$ と $x$ の\R{項}{こう}をくくり、かっこの\R{中}{なか}を\R{平方完成}{へい|ほう|かん|せい}します。\par
\shiki（1）\nobreak $2(x^2-4x)+5=2\{(x-2)^2-4\}+5$\par
\hspace*{2.6\zw}\phantom{（1）}$=2(x-2)^2-8+5=2(x-2)^2-3$\par
\hspace*{2.6\zw}（2）\nobreak $h=-5(t^2-2t)+1=-5\{(t-1)^2-1\}+1$\par
\hspace*{2.6\zw}\phantom{（2）}$=-5(t-1)^2+5+1=-5(t-1)^2+6$\par
\hspace*{2.6\zw}\phantom{（2）}$5(t-1)^2$ は0かプラス。6からそれをひくので、\par
\hspace*{2.6\zw}\phantom{（2）}$h$ は $t=1$ のときにいちばん\R{大}{おお}きく6\par
\kotae（1）\nobreak $2(x-2)^2-3$\qquad（2）\nobreak 1\R{秒後}{びょう|ご}に\qty{6}{m}
\end{reidai}

\begin{jibun}
（1）$3x^2+12x+7$ を\R{平方完成}{へい|ほう|かん|せい}しましょう。\par
（2）\R{高}{たか}さ\qty{2}{m}の\R{所}{ところ}から\qty{10}{m/s}で\R{真上}{ま|うえ}に\R{投}{な}げたボールの\R{高}{たか}さは $h=10t-5t^2+2$（m）です。いちばん\R{高}{たか}くなるのは\R{何秒後}{なん|びょう|ご}で、\R{何}{なん}mですか。
\bkjans{（1）\nobreak $3\{(x+2)^2-4\}+7=3(x+2)^2-5$。（2）\nobreak $h=-5(t-1)^2+7$ で、1\R{秒後}{びょう|ご}に\qty{7}{m}}
\end{jibun}



%% ---------------- まちがえやすいところ ----------------
\bksection{tsumazuki}{まちがえやすいところをチェック}\label{bk:3-8:tsuma}

\begin{tsumazuki}{$(1+h)^2$ のまん\R{中}{なか}をわすれる}
\bkng{$(1+h)^2=1+h^2$}
\bkok{$(1+h)^2=1+2h+h^2$}
$h=0.1$ なら $1.1^2=1.21$ ですが、$1+0.1^2=1.01$ で\R{合}{あ}いません。\R{正方形}{せい|ほう|けい}の\R{図}{ず}（\R{図}{ず}\ref{bk:3-8:fig-sei}）のオレンジの\R{長方形}{ちょう|ほう|けい}2まい $2h$ が\R{抜}{ぬ}けています。
\end{tsumazuki}

\begin{tsumazuki}{$(a+b)(a-b)$ の\R{符号}{ふ|ごう}}
\bkng{$(x+3)(x-3)=x^2+9$}
\bkok{$(x+3)(x-3)=x^2-9$}
\R{最後}{さい|ご}の\R{組}{く}み\R{合}{あ}わせは $3\times(-3)=-9$ です。プラス×マイナスはマイナスです（2-2）。$x=1$ で\R{確}{たし}かめると、$4\times(-2)=-8$、$1-9=-8$ です。
\end{tsumazuki}

\begin{tsumazuki}{「$x=1$ を\R{入}{い}れたから、$x=1$ のときだけの\R{話}{はなし}？」}
\bkng{$x=1$ を\R{入}{い}れて\R{求}{もと}めた $A=3$ は、$x=1$ のときだけ\R{正}{ただ}しい}
\bkok{\R{恒等式}{こう|とう|しき}なので、$A=3$、$B=2$ はどの $x$ でも\R{正}{ただ}しい}
$x=1$ は $B$ を\R{消}{け}すために\R{選}{えら}んだだけ。$x=10$ でも、\R{左}{ひだり}は51、\R{右}{みぎ}は $33+18=51$ です。
\end{tsumazuki}

\begin{tsumazuki}{\R{平方完成}{へい|ほう|かん|せい}で、たした\R{数}{かず}を\R{返}{かえ}しわすれる}
\bkng{$x^2+6x+10=(x+3)^2+10$}
\bkok{$x^2+6x+10=(x+3)^2-9+10=(x+3)^2+1$}
$(x+3)^2$ を\R{展開}{てん|かい}すると $x^2+6x+9$ で、9が\R{多}{おお}すぎます。うめた\R{角}{かど}の9（\R{図}{ず}\ref{bk:3-8:fig-hei}の\R{赤}{あか}）をひいて\R{返}{かえ}します。
\end{tsumazuki}



\begin{butsuri}
\par\addvspace{3pt}\textbf{\R{熱}{ねつ}でふくらむ\R{鉄}{てつ}の\R{板}{いた}}\par\noindent
\begin{minipage}[c]{0.66\linewidth}
\R{鉄}{てつ}は、\R{温度}{おん|ど}が\qty{1}{\degreeCelsius}\R{上}{あ}がるごとに、1 mあたり\R{約}{やく}0.012 mmのびます。1\R{辺}{ぺん}\qty{1000}{mm}の\R{正方形}{せい|ほう|けい}の\R{鉄板}{てっ|ぱん}が\qty{50}{\degreeCelsius}あたたまると、1\R{辺}{ぺん}が\qty{0.6}{mm}のびます（$0.012\times50=0.6$）。ふえる\R{面積}{めん|せき}は、$(a+h)^2-a^2=2ah+h^2$ の $a=1000$、$h=0.6$ で、
\bkthriline{$2\times1000\times0.6+0.6^2=1200+0.36$（mm$^2$）}
です。$h^2$ の0.36は、1200にくらべてとても\R{小}{ちい}さいので、ほとんど $2ah$ だけで\R{計算}{けい|さん}できます。\R{橋}{はし}やレールのつなぎ\R{目}{め}にすき\R{間}{ま}をあけておく\R{大}{おお}きさは、このような\R{計算}{けい|さん}で\R{決}{き}めます（5-8の「\R{小}{ちい}さいあいだは、かけ\R{算}{ざん}で\R{近似}{きん|じ}」につながります）。
\end{minipage}\hfill
\begin{minipage}[c]{0.3\linewidth}\centering
\begin{tikzpicture}[x=1mm, y=1mm]
  \fill[bkBlue!14] (0,4) rectangle (26,30);
  \fill[bkOrange!35] (26,4) rectangle (30,30);
  \fill[bkOrange!35] (0,0) rectangle (26,4);
  \fill[bkRed!45] (26,0) rectangle (30,4);
  \draw[bkInk, line width=0.7pt] (0,0) rectangle (30,30);
  \draw[bkInk, line width=0.4pt] (26,0) -- (26,30) (0,4) -- (30,4);
  \node[font=\footnotesize] at (13,17) {$a^2$};
  \bkbrace{(0,31.5)}{(26,31.5)}{$a$}{1}\bkbrace{(26,31.5)}{(30,31.5)}{$h$}{1}
\end{tikzpicture}
\end{minipage}
\end{butsuri}

%% ---------------- 第13回：部分分数分解（9-7 への橋） ----------------
\bksection{doguu}{\R{分数}{ぶん|すう}を\R{分}{わ}ける：\R{部分分数分解}{ぶ|ぶん|ぶん|すう|ぶん|かい}}

\leavevmode\label{bk:3-8:bubun}\R{道具}{どう|ぐ}7〜9を\R{全部}{ぜん|ぶ}\R{使}{つか}うと、\R{分数}{ぶん|すう}の\R{式}{しき}を、かんたんな\R{分数}{ぶん|すう}の\R{和}{わ}や\R{差}{さ}に\R{分}{わ}けられます。\R{第}{だい}9\R{巻}{かん}で\R{機械}{き|かい}の\R{動}{うご}きを\R{調}{しら}べるときに\R{使}{つか}う\R{計算}{けい|さん}です。

\bkthristep{1}{\R{数}{かず}の\R{分数}{ぶん|すう}で\R{見}{み}る}
\noindent\begin{minipage}[c]{0.55\linewidth}
\R{通分}{つう|ぶん}すると（2-4）、
$\dfrac{1}{2}-\dfrac{1}{3}=\dfrac{3}{6}-\dfrac{2}{6}=\dfrac{1}{6}$ です。6は $2\times3$ なので、$\dfrac{1}{2\times3}=\dfrac{1}{2}-\dfrac{1}{3}$。\R{同}{おな}じように、$\dfrac{1}{3\times4}=\dfrac{1}{3}-\dfrac{1}{4}$ です。\R{分母}{ぶん|ぼ}がかけ\R{算}{ざん}の\R{分数}{ぶん|すう}が、2つの\R{分数}{ぶん|すう}のひき\R{算}{ざん}に\R{分}{わ}けられました。
\end{minipage}\hfill
\begin{minipage}[c]{0.42\linewidth}\centering
\begin{tikzpicture}[x=0.45mm, y=1mm]
  \fill[bkGray!12] (0,14) rectangle (108,18);
  \fill[bkBlue!25] (0,7) rectangle (54,11);
  \fill[bkRed!35] (36,7) rectangle (54,11);
  \fill[bkOrange!30] (0,0) rectangle (36,4);
  \draw[bkInk, line width=0.6pt] (0,14) rectangle (108,18) (0,7) rectangle (54,11) (0,0) rectangle (36,4);
  \draw[bkRed!80!black, line width=0.6pt, dashed] (36,4) -- (36,11);
  \node[font=\footnotesize, anchor=west] at (110,16) {1};
  \node[font=\footnotesize, anchor=west] at (56,9) {$\frac{1}{2}$};
  \node[font=\footnotesize, anchor=west] at (38,2) {$\frac{1}{3}$};
  \node[font=\footnotesize, text=bkRed!80!black, anchor=west] at (62,2.5) {\R{赤}{あか}が $\frac{1}{6}$};
\end{tikzpicture}
\bkfigcaption{$\frac{1}{2}-\frac{1}{3}=\frac{1}{6}$}
\end{minipage}

\bkthristep{2}{\R{文字}{も|じ}の\R{式}{しき}で：わからない\R{数}{かず}を $A$、$B$ とおく}
$\dfrac{1}{(x+1)(x+2)}$ を、\R{分母}{ぶん|ぼ}が $x+1$ と $x+2$ の2つの\R{分数}{ぶん|すう}に\R{分}{わ}けます。\R{分子}{ぶん|し}はわからないので $A$、$B$ とおき、
\bkthriline{$\dfrac{1}{(x+1)(x+2)}=\dfrac{A}{x+1}+\dfrac{B}{x+2}$}
がどんな $x$ でも\R{成}{な}り\R{立}{た}つ（\R{恒等式}{こう|とう|しき}になる）ように $A$、$B$ を\R{決}{き}めます。\R{両辺}{りょう|へん}に $(x+1)(x+2)$ をかけて\R{分母}{ぶん|ぼ}をはらうと（2-14）、\R{左辺}{さ|へん}は1、\R{右辺}{う|へん}は\R{約分}{やく|ぶん}されて $A(x+2)$ と $B(x+1)$ になります。
\bkthriwaku{$1=A(x+2)+B(x+1)$\quad（\R{分数}{ぶん|すう}のない\R{恒等式}{こう|とう|しき}）}

\noindent\textbf{\R{方法}{ほう|ほう}1：\R{係数}{けい|すう}を\R{比}{くら}べる（\R{道具}{どう|ぐ}8）}\quad \R{右辺}{う|へん}$=(A+B)x+(2A+B)$、\R{左辺}{さ|へん}$=0x+1$ なので、$A+B=0$、$2A+B=1$。2つめから1つめをひいて $A=1$（3-6）、だから $B=-1$。

\noindent\textbf{\R{方法}{ほう|ほう}2：\R{数}{かず}を\R{代入}{だい|にゅう}する（\R{道具}{どう|ぐ}9）}\quad $x=-1$ で $B$ の\R{項}{こう}が\R{消}{き}えて $1=A\times1$、$A=1$。$x=-2$ で $A$ の\R{項}{こう}が\R{消}{き}えて $1=B\times(-1)$、$B=-1$（\R{分母}{ぶん|ぼ}をはらった\R{式}{しき}なので、$x=-1$、$-2$ を\R{入}{い}れてよい。\R{道具}{どう|ぐ}9のせいかくに\R{言}{い}うと②）。

\noindent\begin{minipage}[c]{0.5\linewidth}
2\R{通}{とお}りとも $A=1$、$B=-1$ なので、
\[ \frac{1}{(x+1)(x+2)}=\frac{1}{x+1}-\frac{1}{x+2} \]
このように\R{分数}{ぶん|すう}の\R{式}{しき}を\R{分}{わ}けることを\textbf{\R{部分分数分解}{ぶ|ぶん|ぶん|すう|ぶん|かい}}といいます。\R{右}{みぎ}の\R{表}{ひょう}で\R{確}{たし}かめると、$x=1$ の\R{行}{ぎょう}は①そのものです。
\end{minipage}\hfill
\begin{minipage}[c]{0.47\linewidth}\centering
\begin{bktable}{colspec={Q[c,m,wd=6mm] Q[c,m,wd=17mm] Q[c,m,wd=24mm]}, rowsep=0.5pt, cells={font=\small}}
$x$ & \R{左辺}{さ|へん} & \R{右辺}{う|へん} \\
0 & $\frac{1}{1\times2}=\frac{1}{2}$ & $1-\frac{1}{2}=\frac{1}{2}$ \\
1 & $\frac{1}{2\times3}=\frac{1}{6}$ & $\frac{1}{2}-\frac{1}{3}=\frac{1}{6}$ \\
2 & $\frac{1}{3\times4}=\frac{1}{12}$ & $\frac{1}{3}-\frac{1}{4}=\frac{1}{12}$ \\
3 & $\frac{1}{4\times5}=\frac{1}{20}$ & $\frac{1}{4}-\frac{1}{5}=\frac{1}{20}$ \\
\end{bktable}
\end{minipage}

\begin{reidai}[\R{部分分数分解}{ぶ|ぶん|ぶん|すう|ぶん|かい}]
$\dfrac{1}{x(x+1)}=\dfrac{A}{x}+\dfrac{B}{x+1}$ が\R{恒等式}{こう|とう|しき}になるように $A$、$B$ を\R{決}{き}め、$\dfrac{1}{x(x+1)}$ を\R{部分分数分解}{ぶ|ぶん|ぶん|すう|ぶん|かい}しましょう。\par
\kangaekata \R{分母}{ぶん|ぼ}をはらってから\R{数}{かず}を\R{代入}{だい|にゅう}し、\R{係数}{けい|すう}を\R{比}{くら}べて\R{確}{たし}かめます。\par
\shiki \R{両辺}{りょう|へん}に $x(x+1)$ をかける：$1=A(x+1)+Bx$\par
\hspace*{2.6\zw}$x=0$ を\R{入}{い}れる：$1=A\times1+B\times0$ なので $A=1$\par
\hspace*{2.6\zw}$x=-1$ を\R{入}{い}れる：$1=A\times0+B\times(-1)$ なので $B=-1$\par
\hspace*{2.6\zw}\R{確}{たし}かめ：\R{右辺}{う|へん}$=(A+B)x+A$。$A+B=0$、$A=1$ で\R{合}{あ}う\par
\kotae $A=1$、$B=-1$。\quad $\dfrac{1}{x(x+1)}=\dfrac{1}{x}-\dfrac{1}{x+1}$
\end{reidai}

\begin{jibun}
$\dfrac{3}{(x-1)(x+2)}=\dfrac{A}{x-1}+\dfrac{B}{x+2}$ が\R{恒等式}{こう|とう|しき}になるように $A$、$B$ を\R{決}{き}めて\R{部分分数分解}{ぶ|ぶん|ぶん|すう|ぶん|かい}し、\R{係数}{けい|すう}を\R{比}{くら}べる\R{方法}{ほう|ほう}でも\R{確}{たし}かめましょう。
\bkjans{$3=A(x+2)+B(x-1)$。$x=1$ で $3=3A$、$A=1$。$x=-2$ で $3=-3B$、$B=-1$。$\dfrac{3}{(x-1)(x+2)}=\dfrac{1}{x-1}-\dfrac{1}{x+2}$（$A+B=0$、$2A-B=3$ に\R{合}{あ}う）}
\end{jibun}

%% ---------------- 物理・工学 ----------------
\begin{butsuri}
\par\addvspace{3pt}\textbf{\R{制御工学}{せい|ぎょ|こう|がく}：\R{分}{わ}けると、\R{動}{うご}きが\R{読}{よ}める}\par\noindent
\R{機械}{き|かい}を\R{思}{おも}いどおりに\R{動}{うご}かす\R{計算}{けい|さん}（\R{制御}{せい|ぎょ}）では、ラプラス\R{変換}{へん|かん}で、\R{動}{うご}きを $s$ という\R{文字}{も|じ}の\R{分数}{ぶん|すう}の\R{式}{しき}にします（$s$ とラプラス\R{変換}{へん|かん}は\R{第}{だい}9\R{巻}{かん}で\R{学}{まな}びます）。$\dfrac{1}{(s+1)(s+2)}$ のままでは\R{動}{うご}きがわかりませんが、②と\R{同}{おな}じ\R{計算}{けい|さん}で\R{分}{わ}けると、1つずつは\R{表}{ひょう}で\R{引}{ひ}ける\R{形}{かたち}で、「すぐ\R{落}{お}ち\R{着}{つ}く\R{動}{うご}き」（\R{倍々}{ばい|ばい}でへる\R{指数関数}{し|すう|かん|すう}、3-11。5-14の $e$ で $e^{-t}$、$e^{-2t}$ と\R{書}{か}く）にもどせます。\R{全体}{ぜん|たい}は、その\R{差}{さ}です（9-7）。
\begin{bkcenter}
\begin{tikzpicture}[x=1mm, y=1mm]
  \node[bk node, minimum height=11mm, font=\small, inner xsep=1.5mm] (big) at (0,0) {$\dfrac{1}{(s+1)(s+2)}$};
  \node[bk node, minimum height=11mm, font=\small, draw=bkBlue, inner xsep=1.5mm] (p1) at (31,0) {$\dfrac{1}{s+1}$};
  \node[font=\Large] at (40.5,0) {$-$};
  \node[bk node, minimum height=11mm, font=\small, draw=bkGreen, inner xsep=1.5mm] (p2) at (50,0) {$\dfrac{1}{s+2}$};
  \draw[bk link] (big) -- node[bk link label, above, font=\sffamily\scriptsize] {\R{分}{わ}ける} (p1);
  \draw[bk link] (57.5,0) -- node[bk link label, above, font=\sffamily\scriptsize] {\R{表}{ひょう}で} node[bk link label, below, font=\sffamily\scriptsize] {もどす} (69,0);
  \begin{axis}[at={(75.5mm,-17mm)}, anchor=south west, bkplot, width=50mm, height=40mm, xmin=0, xmax=4.4, ymin=0, ymax=1.12, xtick={1,2,3,4}, ytick={0.5,1}, xlabel={$t$}, grid=none]
    \addplot[bkBlue, line width=1.2pt, domain=0:4] {exp(-x)};
    \addplot[bkGreen, line width=1.2pt, domain=0:4] {exp(-2*x)};
    \addplot[bkRed, line width=1.2pt, dashed, domain=0:4] {exp(-x)-exp(-2*x)};
    \draw[bkGreen!50!black, line width=0.4pt] (axis cs:0.4,0.449) -- (axis cs:1.2,0.8);
    \node[font=\sffamily\scriptsize, text=bkGreen!50!black, anchor=west, inner sep=1pt] at (axis cs:1.2,0.8) {$e^{-2t}$};
    \draw[bkBlue!80!black, line width=0.4pt] (axis cs:1.5,0.223) -- (axis cs:2.2,0.55);
    \node[font=\sffamily\scriptsize, text=bkBlue!80!black, anchor=west, inner sep=1pt] at (axis cs:2.2,0.55) {$e^{-t}$};
    \draw[bkRed!80!black, line width=0.4pt] (axis cs:2.4,0.083) -- (axis cs:3.0,0.32);
    \node[font=\sffamily\scriptsize, text=bkRed!80!black, anchor=west, inner sep=1pt] at (axis cs:3.0,0.32) {\R{差}{さ}};
  \end{axis}
\end{tikzpicture}
\end{bkcenter}
\end{butsuri}


%% ---------------- 章末問題 ----------------
\bkendproblems
\bkproblemgroup{\R{基本}{き|ほん}}{1}

\bkproblem \R{次}{つぎ}の\R{計算}{けい|さん}をしましょう（\bkreidairef{1}）。\quad（1）$5a\times(-2ab)$\quad（2）$(-3x)^2$
\bkpans{（1）\nobreak $-10a^2b$。（2）\nobreak $9x^2$}

\bkproblem \R{次}{つぎ}の\R{式}{しき}を\R{展開}{てん|かい}しましょう（\bkreidairef{2}・\bkreidairef{3}）。\par
\hspace*{1\zw}（1）\nobreak $(x+2)(y-5)$\quad（2）\nobreak $(2x-1)(x+3)$\quad（3）\nobreak $(x-5)^2$\par
\hspace*{1\zw}（4）\nobreak $(6-2x)^2$\quad（5）\nobreak $(t+5)(t-3)$\quad（6）\nobreak $(x+9)(x-9)$
\bkpans{（1）\nobreak $xy-5x+2y-10$。（2）\nobreak $2x^2+5x-3$。（3）\nobreak $x^2-10x+25$。（4）\nobreak $36-24x+4x^2$。（5）\nobreak $t^2+2t-15$。（6）\nobreak $x^2-81$}

\bkproblem くふうして\R{計算}{けい|さん}しましょう（\bkreidairef{4}）。\quad（1）$201^2$\quad（2）$1.05^2$\quad（3）$31\times29$
\bkpans{（1）\nobreak $(200+1)^2=40000+400+1=40401$。（2）\nobreak $(1+0.05)^2=1+0.1+0.0025=1.1025$。（3）\nobreak $(30+1)(30-1)=900-1=899$}

\bkproblem \R{共通因数}{きょう|つう|いん|すう}でくくりましょう（\bkreidairef{5}）。\quad（1）$8x^2-24x+16$\quad（2）$3x^2+15x$
\bkpans{（1）\nobreak $8(x^2-3x+2)$。（2）\nobreak $3x(x+5)$}

\bkproblem \R{次}{つぎ}の\R{等式}{とう|しき}は、\R{恒等式}{こう|とう|しき}ですか、\R{方程式}{ほう|てい|しき}ですか。\R{方程式}{ほう|てい|しき}なら、\R{成}{な}り\R{立}{た}つ $x$ も\R{求}{もと}めましょう（\bkreidairef{6}）。\quad（1）$(x+4)(x-4)=x^2-16$\quad（2）$(x+1)^2=x^2+1$
\bkpans{（1）\nobreak \R{恒等式}{こう|とう|しき}。（2）\nobreak $x^2+2x+1=x^2+1$ から $2x=0$ で、$x=0$ のときだけ\R{成}{な}り\R{立}{た}つ\R{方程式}{ほう|てい|しき}}

\bkproblem \R{次}{つぎ}の\R{式}{しき}を\R{平方完成}{へい|ほう|かん|せい}し、いちばん\R{小}{ちい}さい\R{値}{あたい}と、そのときの $x$ を\R{求}{もと}めましょう（\bkreidairef{9}）。\quad（1）$x^2-4x+7$\quad（2）$x^2+2x-3$
\bkpans{（1）\nobreak $(x-2)^2+3$、$x=2$ のとき3。（2）\nobreak $(x+1)^2-4$、$x=-1$ のとき $-4$}

\bkproblemgroup{チャレンジ}{2}

\bkproblem $7x+5=A(x+2)+B(x-1)$ が\R{恒等式}{こう|とう|しき}になるように、$A$、$B$ を\R{決}{き}めましょう。\R{数}{かず}を\R{代入}{だい|にゅう}する\R{方法}{ほう|ほう}と、\R{係数}{けい|すう}を\R{比}{くら}べる\R{方法}{ほう|ほう}の\R{両方}{りょう|ほう}で\R{求}{もと}めましょう（\bkreidairef{7}・\bkreidairef{8}）。
\bkpans{\R{代入}{だい|にゅう}：$x=1$ で $12=3A$、$A=4$。$x=-2$ で $-9=-3B$、$B=3$。\R{係数}{けい|すう}：$(A+B)x+(2A-B)$ から $A+B=7$、$2A-B=5$。たして $3A=12$、$A=4$、$B=3$}

\bkproblem \R{真上}{ま|うえ}に\R{投}{な}げたボールの\R{高}{たか}さが $h=30t-5t^2$（m）のとき、\R{平方完成}{へい|ほう|かん|せい}して、いちばん\R{高}{たか}くなる\R{時刻}{じ|こく}と\R{高}{たか}さを\R{求}{もと}めましょう（\bkreidairef{10}）。
\bkpans{$h=-5(t^2-6t)=-5\{(t-3)^2-9\}=-5(t-3)^2+45$。3\R{秒後}{びょう|ご}に\qty{45}{m}}

\bkproblem $2x^2+12x+13$ を\R{平方完成}{へい|ほう|かん|せい}し、いちばん\R{小}{ちい}さい\R{値}{あたい}を\R{求}{もと}めましょう（\bkreidairef{10}）。
\bkpans{$2(x^2+6x)+13=2\{(x+3)^2-9\}+13=2(x+3)^2-5$。$x=-3$ のとき $-5$}

\begin{furikaeri}
\begin{checklist}
\item かっこのある\R{式}{しき}を\R{展開}{てん|かい}し、\R{面積}{めん|せき}の\R{図}{ず}で\R{説明}{せつ|めい}できる。\bkfurref{bk:3-8:kumiawase}
\item \R{恒等式}{こう|とう|しき}の\R{係数}{けい|すう}を、2\R{通}{とお}りの\R{方法}{ほう|ほう}で\R{決}{き}められる。\bkfurref{bk:3-8:koutou}
\item \R{平方完成}{へい|ほう|かん|せい}して、いちばん\R{小}{ちい}さい・\R{大}{おお}きい\R{値}{あたい}を\R{読}{よ}める。\bkfurref{bk:3-8:heihou}
\end{checklist}
\end{furikaeri}

\bkfukidashi{mina}{\R{式}{しき}の\R{形}{かたち}を\R{変}{か}えると、\R{見}{み}えることが\R{変}{か}わるね。\R{次}{つぎ}の\R{章}{しょう}では、\R{平方完成}{へい|ほう|かん|せい}から\R{放物線}{ほう|ぶつ|せん}のてっぺんを\R{読}{よ}むよ。}
